
\documentclass[UTF8,11pt,a4paper]{ctexart}
\usepackage[a4paper,left=23mm,right=23mm,top=24mm,bottom=24mm]{geometry}
\usepackage{fontspec}
\setmainfont{DejaVu Sans}
\setCJKmainfont{Noto Serif CJK SC}
\usepackage{amsmath,amssymb}
\usepackage{xcolor}
\usepackage[colorlinks=true,linkcolor=blue!45!black,unicode=true]{hyperref}
\usepackage{enumitem}
\usepackage{needspace}
\usepackage{fancyhdr}
\setlength{\parindent}{2em}
\setlength{\parskip}{0.25em}
\setlist[enumerate]{leftmargin=2.6em,itemsep=0.4em,topsep=0.5em}
\pagestyle{fancy}
\fancyhf{}
\fancyhead[L]{\small 三千弱水 · 大学化学高阶综合题库}
\fancyhead[R]{\small 待专家审校稿}
\fancyfoot[C]{\thepage}
\setlength{\headheight}{15pt}
\sloppy

\begin{document}
\begin{center}
{\LARGE\bfseries 大学化学高阶综合题 · 第15辑 (34题)}\\[0.8em]
{\large 题目、分步解析与数值核验}\\[0.5em]
{\small \textcolor{red!65!black}{审校稿：未经化学专家核准；禁止直接视为正式试卷或出版物}}
\end{center}
\vspace{1em}
\noindent 本册从仓库结构化 JSONL 自动生成。增补和修改应只编辑源题目数据，并重新执行统一编译。
\par 机械方程残差检查只能说明指定模型的内部数值一致，不保证模型现实可行性。第六辑另须完成跨辑模型去重审查。
\tableofcontents
\clearpage
\section*{第15辑 · 题目及解答}
\addcontentsline{toc}{section}{第15辑}
\subsection*{无机化学}
\addcontentsline{toc}{subsection}{第15辑 / 无机化学}
\Needspace{7\baselineskip}
\subsubsection*{UC-ADV15-I226 · LiF的Born–Haber焓循环与晶格形成能}
\noindent\textbf{题目}\par
由Li(s)+½F₂(g)→LiF(s)标准形成焓−617.0 kJ/mol，利用Li升华焓+159.0 kJ/mol、F₂键解离焓+158.0 kJ/mol分子、Li电离能+520.0 kJ/mol、F电子亲和反应焓−328.0 kJ/mol建立完整Born–Haber循环。求形成晶格时的放热焓与相反的晶格分离焓、非晶格步骤总和、½F₂原子化焓、晶格对形成反应焓的贡献比例以及将电子亲和放热误记为吸热时造成的数值偏差。所有数值以一摩尔LiF式量为基准，并忽略温度热容校正。\par
\medskip\noindent\textbf{分步解答}
\begin{enumerate}
\item 题目输入与理想假设：由Li(s)+½F₂(g)→LiF(s)标准形成焓−617.0 kJ/mol，利用Li升华焓+159.0 kJ/mol、F₂键解离焓+158.0 kJ/mol分子、Li电离能+520.0 kJ/mol、F电子亲和反应焓−328.0 kJ/mol建立完整Born–Haber循环。求形成晶格时的放热焓与相反的晶格分离焓、非晶格步骤总和、½F₂原子化焓、晶格对形成反应焓的贡献比例以及将电子亲和放热误记为吸热时造成的数值偏差。所有数值以一摩尔LiF式量为基准，并忽略温度热容校正。
\item 根据电荷、能量或物料守恒建立方程：ΔHf=ΔHsub+½D+IE+EA+Ulattice，故Ulattice=−617−(159+79+520−328)。
\item 统一采用题中给定标准态并把分数、浓度与活动度写为无量纲反应商。
\item 将独立变量代入基本方程，解出第一组量：atomization\_half\_F2 = 79.000000000000 kJ/mol；Li\_sublimation = 159.00000000000 kJ/mol；Li\_ionization = 520.00000000000 kJ/mol。
\item 计算平衡或结构中间量：F\_electron\_affinity = -328.00000000000 kJ/mol；nonlattice\_sum = 430.00000000000 kJ/mol；lattice\_formation = -1047.0000000000 kJ/mol。
\item 进一步计算其余可观测量：lattice\_separation = 1047.0000000000 kJ/mol；formation\_enthalpy = -617.00000000000 kJ/mol；absolute\_ratio\_lattice\_formation = 1.6969205834684 ratio；wrong\_sign\_bias = -656.00000000000 kJ/mol。
\item 把计算值分别代回原始离子、电荷、分配与能量约束检验守恒和符号。
\item 通过温度、电子数或配体含量的极限检查模型自洽性。
\item 排除题目未规定的配位异构、溶剂化、活度修正和相变动力学，并指出实际实验需要补充的测定。
\item 本题需专家复查科学性和跨辑独立性，数值脚本运行不等于专家批准。
\end{enumerate}
\noindent\textbf{答案摘要}\par
atomization\_half\_F2 = 79 kJ/mol；Li\_sublimation = 159 kJ/mol；Li\_ionization = 520 kJ/mol；F\_electron\_affinity = -328 kJ/mol；nonlattice\_sum = 430 kJ/mol；lattice\_formation = -1047 kJ/mol；lattice\_separation = 1047 kJ/mol；formation\_enthalpy = -617 kJ/mol；absolute\_ratio\_lattice\_formation = 1.69692058 ratio；wrong\_sign\_bias = -656 kJ/mol\par
\noindent\textbf{模型前提：}
仅使用理想模型和给定常数，专家审核未完成。\par
\noindent\textcolor{red!65!black}{[跨辑去重未完成 · 草稿]} \par
\medskip\hrule\medskip
\Needspace{7\baselineskip}
\subsubsection*{UC-ADV15-I227 · Fe(II)热自旋交叉的高低自旋态分数}
\noindent\textbf{题目}\par
某六配位Fe(II)配合物有低自旋LS与高自旋HS两热力学状态，HS相对LS有ΔH=+20.00 kJ/mol和ΔS=+60.00 J/(mol K)，两态简并度已包含在ΔS中，不额外乘5。T=250.0 K与350.0 K，定义K=[HS]/[LS]=exp[−(ΔH−TΔS)/(RT)]。求两个温度的自由能差、K、HS/LS分率以及两温度平衡常数比、HS分率差和K=1时的交叉温度。需解释高温促进高熵态而忽略协同自旋相变、晶格滞后与动力学限制。\par
\medskip\noindent\textbf{分步解答}
\begin{enumerate}
\item 题目输入与理想假设：某六配位Fe(II)配合物有低自旋LS与高自旋HS两热力学状态，HS相对LS有ΔH=+20.00 kJ/mol和ΔS=+60.00 J/(mol K)，两态简并度已包含在ΔS中，不额外乘5。T=250.0 K与350.0 K，定义K=[HS]/[LS]=exp[−(ΔH−TΔS)/(RT)]。求两个温度的自由能差、K、HS/LS分率以及两温度平衡常数比、HS分率差和K=1时的交叉温度。需解释高温促进高熵态而忽略协同自旋相变、晶格滞后与动力学限制。
\item 根据电荷、能量或物料守恒建立方程：ΔG=ΔH−TΔS，fHS=K/(1+K)，Tcross=ΔH/ΔS。
\item 统一采用题中给定标准态并把分数、浓度与活动度写为无量纲反应商。
\item 将独立变量代入基本方程，解出第一组量：DG\_250 = 5.0000000000000 kJ/mol；DG\_350 = -1.0000000000000 kJ/mol；K\_250 = 0.090225146846218 ratio。
\item 计算平衡或结构中间量：K\_350 = 1.4100642914567 ratio；fraction\_HS\_250 = 0.082758269800710 fraction；fraction\_HS\_350 = 0.58507330964372 fraction。
\item 进一步计算其余可观测量：fraction\_LS\_250 = 0.91724173019929 fraction；ratio\_K\_350\_250 = 15.628284804680 ratio；cross\_temperature = 333.33333333333 K；increase\_HS\_fraction = 0.50231503984301 fraction。
\item 把计算值分别代回原始离子、电荷、分配与能量约束检验守恒和符号。
\item 通过温度、电子数或配体含量的极限检查模型自洽性。
\item 排除题目未规定的配位异构、溶剂化、活度修正和相变动力学，并指出实际实验需要补充的测定。
\item 本题需专家复查科学性和跨辑独立性，数值脚本运行不等于专家批准。
\end{enumerate}
\noindent\textbf{答案摘要}\par
DG\_250 = 5 kJ/mol；DG\_350 = -1 kJ/mol；K\_250 = 0.0902251468 ratio；K\_350 = 1.41006429 ratio；fraction\_HS\_250 = 0.0827582698 fraction；fraction\_HS\_350 = 0.58507331 fraction；fraction\_LS\_250 = 0.91724173 fraction；ratio\_K\_350\_250 = 15.6282848 ratio；cross\_temperature = 333.333333 K；increase\_HS\_fraction = 0.50231504 fraction\par
\noindent\textbf{模型前提：}
仅使用理想模型和给定常数，专家审核未完成。\par
\noindent\textcolor{red!65!black}{[跨辑去重未完成 · 草稿]} \par
\medskip\hrule\medskip
\Needspace{7\baselineskip}
\subsubsection*{UC-ADV15-I228 · Cu(I)歧化反应的标准平衡常数与浓度驱动力}
\noindent\textbf{题目}\par
在298.15 K考察2Cu⁺(aq)⇌Cu²⁺(aq)+Cu(s)，还原电势Cu²⁺/Cu⁺为0.1530 V，Cu⁺/Cu为0.5210 V。理想水溶液给定自由[Cu⁺]=0.001000 M、[Cu²⁺]=0.01000 M，铜固体活动度取1。求歧化标准电动势E°、lnK、log10K、K、当前反应商Q、Nernst实际电动势、每mol反应的实际及标准Gibbs自由能、若Cu²⁺被固定在原值而达到平衡所需的Cu⁺浓度及实际歧化自发方向标志。正确使用总反应n=1。\par
\medskip\noindent\textbf{分步解答}
\begin{enumerate}
\item 题目输入与理想假设：在298.15 K考察2Cu⁺(aq)⇌Cu²⁺(aq)+Cu(s)，还原电势Cu²⁺/Cu⁺为0.1530 V，Cu⁺/Cu为0.5210 V。理想水溶液给定自由[Cu⁺]=0.001000 M、[Cu²⁺]=0.01000 M，铜固体活动度取1。求歧化标准电动势E°、lnK、log10K、K、当前反应商Q、Nernst实际电动势、每mol反应的实际及标准Gibbs自由能、若Cu²⁺被固定在原值而达到平衡所需的Cu⁺浓度及实际歧化自发方向标志。正确使用总反应n=1。
\item 根据电荷、能量或物料守恒建立方程：E°=E°Cu+/Cu−E°Cu2+/Cu+=0.368；K=exp(FE°/RT)，Q=[Cu²⁺]/[Cu⁺]²，E=E°−RT/F ln Q。
\item 统一采用题中给定标准态并把分数、浓度与活动度写为无量纲反应商。
\item 将独立变量代入基本方程，解出第一组量：E\_standard = 0.36800000000000 V；ln\_K = 14.323201974384 log；log10\_K = 6.2204875806608 log10。
\item 计算平衡或结构中间量：K\_equilibrium = 1661451.1652573 ratio；Q\_reaction = 10000.000000000 ratio；E\_actual = 0.13136260125711 V。
\item 进一步计算其余可观测量：DG\_actual\_kJ = -12.674564210440 kJ/mol；DG\_standard\_kJ = -35.506602220160 kJ/mol；CuI\_eq\_M = 0.000077581149415667 mol/L；spontaneous\_flag = 1.0000000000000 boolean。
\item 把计算值分别代回原始离子、电荷、分配与能量约束检验守恒和符号。
\item 通过温度、电子数或配体含量的极限检查模型自洽性。
\item 排除题目未规定的配位异构、溶剂化、活度修正和相变动力学，并指出实际实验需要补充的测定。
\item 本题需专家复查科学性和跨辑独立性，数值脚本运行不等于专家批准。
\end{enumerate}
\noindent\textbf{答案摘要}\par
E\_standard = 0.368 V；ln\_K = 14.323202 log；log10\_K = 6.22048758 log10；K\_equilibrium = 1661451.17 ratio；Q\_reaction = 10000 ratio；E\_actual = 0.131362601 V；DG\_actual\_kJ = -12.6745642 kJ/mol；DG\_standard\_kJ = -35.5066022 kJ/mol；CuI\_eq\_M = 7.75811494e-05 mol/L；spontaneous\_flag = 1 boolean\par
\noindent\textbf{模型前提：}
仅使用理想模型和给定常数，专家审核未完成。\par
\noindent\textcolor{red!65!black}{[跨辑去重未完成 · 草稿]} \par
\medskip\hrule\medskip
\Needspace{7\baselineskip}
\subsubsection*{UC-ADV15-I229 · 一价与二价离子的Born溶剂化自由能比较}
\noindent\textbf{题目}\par
在水连续介质近似中取相对介电常数εr=78.50，Born溶剂化自由能ΔG=−NA z²e²(1−1/εr)/(8π ε0 r)，离子球半径分别是一价r1=0.1000 nm、二价r2=0.08000 nm。已知NA=6.02214076×10²³ mol⁻¹、e=1.602176634×10⁻¹⁹ C、ε0=8.8541878128×10⁻¹² F/m。计算两离子溶剂化能(kJ/mol)、绝对能量比、各自z²/r(nm⁻¹)、溶剂化能差、同半径二价离子溶剂化能、无限大介电常数极限一价离子能量和有限/无限介电强度比，解释为什么连续模型不包含离子配位水结构及化学作用。\par
\medskip\noindent\textbf{分步解答}
\begin{enumerate}
\item 题目输入与理想假设：在水连续介质近似中取相对介电常数εr=78.50，Born溶剂化自由能ΔG=−NA z²e²(1−1/εr)/(8π ε0 r)，离子球半径分别是一价r1=0.1000 nm、二价r2=0.08000 nm。已知NA=6.02214076×10²³ mol⁻¹、e=1.602176634×10⁻¹⁹ C、ε0=8.8541878128×10⁻¹² F/m。计算两离子溶剂化能(kJ/mol)、绝对能量比、各自z²/r(nm⁻¹)、溶剂化能差、同半径二价离子溶剂化能、无限大介电常数极限一价离子能量和有限/无限介电强度比，解释为什么连续模型不包含离子配位水结构及化学作用。
\item 根据电荷、能量或物料守恒建立方程：Born系数C=NA e²/(8π ε0)；ΔG=−C(z²/r)(1−1/εr)。
\item 统一采用题中给定标准态并把分数、浓度与活动度写为无量纲反应商。
\item 将独立变量代入基本方程，解出第一组量：G\_monovalent\_kJ = -685.82789601526 kJ/mol；G\_divalent\_kJ = -3429.1394800763 kJ/mol；absolute\_ratio = 5.0000000000000 ratio。
\item 计算平衡或结构中间量：z2\_over\_r\_monovalent = 10.000000000000 nm\textasciicircum{}-1；z2\_over\_r\_divalent = 50.000000000000 nm\textasciicircum{}-1；energy\_difference\_kJ = -2743.3115840610 kJ/mol。
\item 进一步计算其余可观测量：G\_divalent\_same\_radius\_kJ = -2743.3115840610 kJ/mol；same\_radius\_energy\_ratio = 4.0000000000000 ratio；G\_monovalent\_eps\_inf\_kJ = -694.67728822191 kJ/mol；finite\_eps\_ratio = 0.98726114649682 ratio。
\item 把计算值分别代回原始离子、电荷、分配与能量约束检验守恒和符号。
\item 通过温度、电子数或配体含量的极限检查模型自洽性。
\item 排除题目未规定的配位异构、溶剂化、活度修正和相变动力学，并指出实际实验需要补充的测定。
\item 本题需专家复查科学性和跨辑独立性，数值脚本运行不等于专家批准。
\end{enumerate}
\noindent\textbf{答案摘要}\par
G\_monovalent\_kJ = -685.827896 kJ/mol；G\_divalent\_kJ = -3429.13948 kJ/mol；absolute\_ratio = 5 ratio；z2\_over\_r\_monovalent = 10 nm\textasciicircum{}-1；z2\_over\_r\_divalent = 50 nm\textasciicircum{}-1；energy\_difference\_kJ = -2743.31158 kJ/mol；G\_divalent\_same\_radius\_kJ = -2743.31158 kJ/mol；same\_radius\_energy\_ratio = 4 ratio；G\_monovalent\_eps\_inf\_kJ = -694.677288 kJ/mol；finite\_eps\_ratio = 0.987261146 ratio\par
\noindent\textbf{模型前提：}
仅使用理想模型和给定常数，专家审核未完成。\par
\noindent\textcolor{red!65!black}{[跨辑去重未完成 · 草稿]} \par
\medskip\hrule\medskip
\Needspace{7\baselineskip}
\subsubsection*{UC-ADV15-I230 · BaSO₄与SrSO₄竞争沉淀的选择性窗口}
\noindent\textbf{题目}\par
在25°C理想水溶液中自由Ba²⁺由外池维持0.005000 M，自由Sr²⁺由外池维持0.002000 M，硫酸根逐渐加入、浓度可调。取BaSO₄的Ksp=1.000×10⁻¹⁰，SrSO₄的Ksp=3.000×10⁻⁷（均按1 M标准态）。求两种沉淀恰开始析出时自由[SO₄²⁻]阈值、阈值比、在自由SO₄=1.000×10⁻⁶ M时两盐离子积及过饱和度，能够只析出BaSO₄的硫酸根浓度窗口上下界及宽度、优先析出盐的判断。说明维持自由金属浓度需要外部化学储库，本题不推算封闭样品总沉淀质量。\par
\medskip\noindent\textbf{分步解答}
\begin{enumerate}
\item 题目输入与理想假设：在25°C理想水溶液中自由Ba²⁺由外池维持0.005000 M，自由Sr²⁺由外池维持0.002000 M，硫酸根逐渐加入、浓度可调。取BaSO₄的Ksp=1.000×10⁻¹⁰，SrSO₄的Ksp=3.000×10⁻⁷（均按1 M标准态）。求两种沉淀恰开始析出时自由[SO₄²⁻]阈值、阈值比、在自由SO₄=1.000×10⁻⁶ M时两盐离子积及过饱和度，能够只析出BaSO₄的硫酸根浓度窗口上下界及宽度、优先析出盐的判断。说明维持自由金属浓度需要外部化学储库，本题不推算封闭样品总沉淀质量。
\item 根据电荷、能量或物料守恒建立方程：阈值硫酸根Ksp/[M²⁺]，Q=[M²⁺][SO₄²⁻]，Q/Ksp>1时过饱和。
\item 统一采用题中给定标准态并把分数、浓度与活动度写为无量纲反应商。
\item 将独立变量代入基本方程，解出第一组量：Ba\_threshold\_M = 2.0000000000000e-8 mol/L；Sr\_threshold\_M = 0.00015000000000000 mol/L；threshold\_ratio\_Sr\_Ba = 7500.0000000000 ratio。
\item 计算平衡或结构中间量：Q\_Ba\_at1eMinus6 = 5.0000000000000e-9 M2；Q\_Sr\_at1eMinus6 = 2.0000000000000e-9 M2；saturation\_Ba = 50.000000000000 ratio。
\item 进一步计算其余可观测量：saturation\_Sr = 0.0066666666666667 ratio；selective\_window\_lower = 2.0000000000000e-8 mol/L；selective\_window\_upper = 0.00015000000000000 mol/L；window\_width = 0.00014998000000000 mol/L。
\item 把计算值分别代回原始离子、电荷、分配与能量约束检验守恒和符号。
\item 通过温度、电子数或配体含量的极限检查模型自洽性。
\item 排除题目未规定的配位异构、溶剂化、活度修正和相变动力学，并指出实际实验需要补充的测定。
\item 本题需专家复查科学性和跨辑独立性，数值脚本运行不等于专家批准。
\end{enumerate}
\noindent\textbf{答案摘要}\par
Ba\_threshold\_M = 2e-08 mol/L；Sr\_threshold\_M = 0.00015 mol/L；threshold\_ratio\_Sr\_Ba = 7500 ratio；Q\_Ba\_at1eMinus6 = 5e-09 M2；Q\_Sr\_at1eMinus6 = 2e-09 M2；saturation\_Ba = 50 ratio；saturation\_Sr = 0.00666666667 ratio；selective\_window\_lower = 2e-08 mol/L；selective\_window\_upper = 0.00015 mol/L；window\_width = 0.00014998 mol/L\par
\noindent\textbf{模型前提：}
仅使用理想模型和给定常数，专家审核未完成。\par
\noindent\textcolor{red!65!black}{[跨辑去重未完成 · 草稿]} \par
\medskip\hrule\medskip
\Needspace{7\baselineskip}
\subsubsection*{UC-ADV15-I231 · ABO₃氧空位与B位还原的电荷补偿}
\noindent\textbf{题目}\par
氧化物ABO₃的B位初始为B³⁺，氧位O²⁻。在还原氛围下化学式变成ABO₃₋δ，δ=0.01000，氧空位按每个产生两电子，每个电子将一个B³⁺还原为B²⁺。忽略其它载流子和缺陷。计算每式量氧空位数、电子数、B²⁺与B³⁺摩尔分率、B平均氧化态、实际氧原子数、氧位空位分数、1 mol式量对应失氧量与电子摩尔数，以及有效空位电荷与补偿电子的净电荷残差。说明这里是电荷守恒示例，实际B位可能发生非局域混合价与电子导电。\par
\medskip\noindent\textbf{分步解答}
\begin{enumerate}
\item 题目输入与理想假设：氧化物ABO₃的B位初始为B³⁺，氧位O²⁻。在还原氛围下化学式变成ABO₃₋δ，δ=0.01000，氧空位按每个产生两电子，每个电子将一个B³⁺还原为B²⁺。忽略其它载流子和缺陷。计算每式量氧空位数、电子数、B²⁺与B³⁺摩尔分率、B平均氧化态、实际氧原子数、氧位空位分数、1 mol式量对应失氧量与电子摩尔数，以及有效空位电荷与补偿电子的净电荷残差。说明这里是电荷守恒示例，实际B位可能发生非局域混合价与电子导电。
\item 根据电荷、能量或物料守恒建立方程：n(e′)=2δ，xBII=2δ，vB=3−2δ，nO=3−δ，电荷补偿2δ−n(e′)=0。
\item 统一采用题中给定标准态并把分数、浓度与活动度写为无量纲反应商。
\item 将独立变量代入基本方程，解出第一组量：oxygen\_vacancies = 0.010000000000000 per\_formula；electron\_count = 0.020000000000000 electrons；BII\_fraction = 0.020000000000000 fraction。
\item 计算平衡或结构中间量：BIII\_fraction = 0.98000000000000 fraction；mean\_B\_oxidation = 2.9800000000000 oxidation；oxygen\_count = 2.9900000000000 atoms。
\item 进一步计算其余可观测量：vacancy\_site\_fraction = 0.0033333333333333 fraction；oxygen\_removed\_mol = 0.010000000000000 mol/mol；electron\_moles = 0.020000000000000 mol/mol；net\_defect\_charge = 0.0000000000000 charge。
\item 把计算值分别代回原始离子、电荷、分配与能量约束检验守恒和符号。
\item 通过温度、电子数或配体含量的极限检查模型自洽性。
\item 排除题目未规定的配位异构、溶剂化、活度修正和相变动力学，并指出实际实验需要补充的测定。
\item 本题需专家复查科学性和跨辑独立性，数值脚本运行不等于专家批准。
\end{enumerate}
\noindent\textbf{答案摘要}\par
oxygen\_vacancies = 0.01 per\_formula；electron\_count = 0.02 electrons；BII\_fraction = 0.02 fraction；BIII\_fraction = 0.98 fraction；mean\_B\_oxidation = 2.98 oxidation；oxygen\_count = 2.99 atoms；vacancy\_site\_fraction = 0.00333333333 fraction；oxygen\_removed\_mol = 0.01 mol/mol；electron\_moles = 0.02 mol/mol；net\_defect\_charge = 0 charge\par
\noindent\textbf{模型前提：}
仅使用理想模型和给定常数，专家审核未完成。\par
\noindent\textcolor{red!65!black}{[跨辑去重未完成 · 草稿]} \par
\medskip\hrule\medskip
\Needspace{7\baselineskip}
\subsubsection*{UC-ADV15-I232 · 配体交换的Eyring速率与活化自由能}
\noindent\textbf{题目}\par
某八面体金属配合物的一阶配体交换反应满足过渡态理论k=(kBT/h)exp(ΔS‡/R)exp(−ΔH‡/(RT))，活化焓+65.00 kJ/mol，活化熵−30.00 J/(mol K)，T分别298.15和350.0 K，kB=1.380649×10⁻²³ J/K，h=6.62607015×10⁻³⁴ J·s。求两温度活化Gibbs自由能、两个一级速度常数与半衰期、速率比、Eyring斜率−ΔH‡/R和过渡态理论前因子比、活化熵的符号。假设无机配位化学机理和活化热力学参数在温区内保持不变。\par
\medskip\noindent\textbf{分步解答}
\begin{enumerate}
\item 题目输入与理想假设：某八面体金属配合物的一阶配体交换反应满足过渡态理论k=(kBT/h)exp(ΔS‡/R)exp(−ΔH‡/(RT))，活化焓+65.00 kJ/mol，活化熵−30.00 J/(mol K)，T分别298.15和350.0 K，kB=1.380649×10⁻²³ J/K，h=6.62607015×10⁻³⁴ J·s。求两温度活化Gibbs自由能、两个一级速度常数与半衰期、速率比、Eyring斜率−ΔH‡/R和过渡态理论前因子比、活化熵的符号。假设无机配位化学机理和活化热力学参数在温区内保持不变。
\item 根据电荷、能量或物料守恒建立方程：ΔG‡=ΔH‡−TΔS‡，k=kBT/h·exp(−ΔG‡/RT)。
\item 统一采用题中给定标准态并把分数、浓度与活动度写为无量纲反应商。
\item 将独立变量代入基本方程，解出第一组量：G298\_kJ = 73.944500000000 kJ/mol；G350\_kJ = 75.500000000000 kJ/mol；rate298 = 0.68983645579233 s\textasciicircum{}-1。
\item 计算平衡或结构中间量：rate350 = 39.387393583370 s\textasciicircum{}-1；half298 = 1.0047992893675 s；half350 = 0.017598198750897 s。
\item 进一步计算其余可观测量：rate\_ratio = 57.096712202794 ratio；Eyring\_slope = -7817.7030779213 K；prefactor\_ratio = 1.1739057521382 ratio；activation\_entropy = -30.000000000000 J/(molK)。
\item 把计算值分别代回原始离子、电荷、分配与能量约束检验守恒和符号。
\item 通过温度、电子数或配体含量的极限检查模型自洽性。
\item 排除题目未规定的配位异构、溶剂化、活度修正和相变动力学，并指出实际实验需要补充的测定。
\item 本题需专家复查科学性和跨辑独立性，数值脚本运行不等于专家批准。
\end{enumerate}
\noindent\textbf{答案摘要}\par
G298\_kJ = 73.9445 kJ/mol；G350\_kJ = 75.5 kJ/mol；rate298 = 0.689836456 s\textasciicircum{}-1；rate350 = 39.3873936 s\textasciicircum{}-1；half298 = 1.00479929 s；half350 = 0.0175981988 s；rate\_ratio = 57.0967122 ratio；Eyring\_slope = -7817.70308 K；prefactor\_ratio = 1.17390575 ratio；activation\_entropy = -30 J/(molK)\par
\noindent\textbf{模型前提：}
仅使用理想模型和给定常数，专家审核未完成。\par
\noindent\textcolor{red!65!black}{[跨辑去重未完成 · 草稿]} \par
\medskip\hrule\medskip
\Needspace{7\baselineskip}
\subsubsection*{UC-ADV15-I233 · 三价金属与两种单齿配体的竞争形成分率}
\noindent\textbf{题目}\par
理想水溶液中金属M³⁺能分别与自由配体A、B形成MA³⁺和MB³⁺，1:1条件形成常数分别βA=1000、βB=300.0 M⁻¹，两个自由配体浓度由大储液池维持[A]=0.002000、[B]=0.004000 M。总金属浓度CT=0.01000 M，温度固定，不发生金属水解、三元络合物和金属沉淀。求两组配位权重、配分函数、三种含金属物种浓度与分率、MA/MB浓度比、金属守恒核验以及全部络合金属占比。区别游离配体和投料配体的浓度。\par
\medskip\noindent\textbf{分步解答}
\begin{enumerate}
\item 题目输入与理想假设：理想水溶液中金属M³⁺能分别与自由配体A、B形成MA³⁺和MB³⁺，1:1条件形成常数分别βA=1000、βB=300.0 M⁻¹，两个自由配体浓度由大储液池维持[A]=0.002000、[B]=0.004000 M。总金属浓度CT=0.01000 M，温度固定，不发生金属水解、三元络合物和金属沉淀。求两组配位权重、配分函数、三种含金属物种浓度与分率、MA/MB浓度比、金属守恒核验以及全部络合金属占比。区别游离配体和投料配体的浓度。
\item 根据电荷、能量或物料守恒建立方程：wA=βA[A]，wB=βB[B]，D=1+wA+wB，物种分率w\_i/D。
\item 统一采用题中给定标准态并把分数、浓度与活动度写为无量纲反应商。
\item 将独立变量代入基本方程，解出第一组量：weight\_A = 2.0000000000000 ratio；weight\_B = 1.2000000000000 ratio；partition\_function = 4.2000000000000 ratio。
\item 计算平衡或结构中间量：metal\_free\_M = 0.0023809523809524 mol/L；MA\_M = 0.0047619047619048 mol/L；MB\_M = 0.0028571428571429 mol/L。
\item 进一步计算其余可观测量：fraction\_free = 0.23809523809524 fraction；fraction\_MA = 0.47619047619048 fraction；fraction\_MB = 0.28571428571429 fraction；MA\_MB\_ratio = 1.6666666666667 ratio；total\_complexed\_fraction = 0.76190476190476 fraction。
\item 把计算值分别代回原始离子、电荷、分配与能量约束检验守恒和符号。
\item 通过温度、电子数或配体含量的极限检查模型自洽性。
\item 排除题目未规定的配位异构、溶剂化、活度修正和相变动力学，并指出实际实验需要补充的测定。
\item 本题需专家复查科学性和跨辑独立性，数值脚本运行不等于专家批准。
\end{enumerate}
\noindent\textbf{答案摘要}\par
weight\_A = 2 ratio；weight\_B = 1.2 ratio；partition\_function = 4.2 ratio；metal\_free\_M = 0.00238095238 mol/L；MA\_M = 0.00476190476 mol/L；MB\_M = 0.00285714286 mol/L；fraction\_free = 0.238095238 fraction；fraction\_MA = 0.476190476 fraction；fraction\_MB = 0.285714286 fraction；MA\_MB\_ratio = 1.66666667 ratio；total\_complexed\_fraction = 0.761904762 fraction\par
\noindent\textbf{模型前提：}
仅使用理想模型和给定常数，专家审核未完成。\par
\noindent\textcolor{red!65!black}{[跨辑去重未完成 · 草稿]} \par
\medskip\hrule\medskip
\Needspace{7\baselineskip}
\subsubsection*{UC-ADV15-I234 · 二价与三价钒离子的标准Nernst氧化还原电位}
\noindent\textbf{题目}\par
298.15 K假想钒离子V³⁺+e⁻⇌V²⁺的标准还原电势E°=−0.2550 V。两种状态下水溶液自由V³⁺与V²⁺活动度分别为(aIII,aII)=(0.02000,0.005000)和(0.002000,0.01000)，按1 M标准态定义无量纲活动度，不存在水解、配位及电极极化。求两个反应商Q=aII/aIII、两个平衡电势、两电位差、两状态氧化态分率、Nernst十倍活动度比斜率RTln10/F、第一状态V(III)/V(II)比与等活动度状态的电势。强调n=1。\par
\medskip\noindent\textbf{分步解答}
\begin{enumerate}
\item 题目输入与理想假设：298.15 K假想钒离子V³⁺+e⁻⇌V²⁺的标准还原电势E°=−0.2550 V。两种状态下水溶液自由V³⁺与V²⁺活动度分别为(aIII,aII)=(0.02000,0.005000)和(0.002000,0.01000)，按1 M标准态定义无量纲活动度，不存在水解、配位及电极极化。求两个反应商Q=aII/aIII、两个平衡电势、两电位差、两状态氧化态分率、Nernst十倍活动度比斜率RTln10/F、第一状态V(III)/V(II)比与等活动度状态的电势。强调n=1。
\item 根据电荷、能量或物料守恒建立方程：E=E°−(RT/F)ln(aII/aIII)，状态分率aIII/(aIII+aII)。
\item 统一采用题中给定标准态并把分数、浓度与活动度写为无量纲反应商。
\item 将独立变量代入基本方程，解出第一组量：Q\_state1 = 0.25000000000000 ratio；Q\_state2 = 5.0000000000000 ratio；E\_state1\_V = -0.21938252244125 V。
\item 计算平衡或结构中间量：E\_state2\_V = -0.29635061090634 V；E2\_minus\_E1 = -0.076968088465098 V；V3\_fraction\_state1 = 0.80000000000000 fraction。
\item 进一步计算其余可观测量：V3\_fraction\_state2 = 0.16666666666667 fraction；slope\_per\_decade\_V = 0.059159349685722 V；V3\_V2\_ratio\_state1 = 4.0000000000000 ratio；E\_equal\_activities = -0.25500000000000 V。
\item 把计算值分别代回原始离子、电荷、分配与能量约束检验守恒和符号。
\item 通过温度、电子数或配体含量的极限检查模型自洽性。
\item 排除题目未规定的配位异构、溶剂化、活度修正和相变动力学，并指出实际实验需要补充的测定。
\item 本题需专家复查科学性和跨辑独立性，数值脚本运行不等于专家批准。
\end{enumerate}
\noindent\textbf{答案摘要}\par
Q\_state1 = 0.25 ratio；Q\_state2 = 5 ratio；E\_state1\_V = -0.219382522 V；E\_state2\_V = -0.296350611 V；E2\_minus\_E1 = -0.0769680885 V；V3\_fraction\_state1 = 0.8 fraction；V3\_fraction\_state2 = 0.166666667 fraction；slope\_per\_decade\_V = 0.0591593497 V；V3\_V2\_ratio\_state1 = 4 ratio；E\_equal\_activities = -0.255 V\par
\noindent\textbf{模型前提：}
仅使用理想模型和给定常数，专家审核未完成。\par
\noindent\textcolor{red!65!black}{[跨辑去重未完成 · 草稿]} \par
\medskip\hrule\medskip
\subsection*{分析化学}
\addcontentsline{toc}{subsection}{第15辑 / 分析化学}
\Needspace{7\baselineskip}
\subsubsection*{UC-ADV15-A226 · 同位素稀释质谱的重原子丰度精确反演}
\noindent\textbf{题目}\par
某单重同位素示踪分子的天然样品中重核原子分数fX=0.01000，实验向样品加入0.002000 mol化学性质相同、重核富集分率fS=0.9000的标记内标。混合物的重原子分数测得fm=0.2500，假设每分子恰有一个被监测原子位点且无同位素分馏、检测器质量歧视和试剂损失。通过总原子数守恒反演原样未知物质的量nX，求两来源重核摩尔数、混合总核摩尔数、混合重核和轻核总量、原样/内标摩尔比、混合重轻原子比和混合分率复算。说明不能直接用重/轻比替代重原子分数。\par
\medskip\noindent\textbf{分步解答}
\begin{enumerate}
\item 实验测量模型和给定条件：某单重同位素示踪分子的天然样品中重核原子分数fX=0.01000，实验向样品加入0.002000 mol化学性质相同、重核富集分率fS=0.9000的标记内标。混合物的重原子分数测得fm=0.2500，假设每分子恰有一个被监测原子位点且无同位素分馏、检测器质量歧视和试剂损失。通过总原子数守恒反演原样未知物质的量nX，求两来源重核摩尔数、混合总核摩尔数、混合重核和轻核总量、原样/内标摩尔比、混合重轻原子比和混合分率复算。说明不能直接用重/轻比替代重原子分数。
\item 完整反演或计量方程：fm=(nX fX+nS fS)/(nX+nS)，nX=nS(fS−fm)/(fm−fX)，nheavy=fm(nX+nS)。
\item 确认原样、标准品和空白的质量或体积关系，所有光谱响应和电位基线按题面规定处理。
\item 第一批直接求得：sample\_mol = 0.005416666666667 mol；spike\_mol = 0.002000000000000 mol；sample\_spike\_ratio = 2.708333333333 ratio。
\item 第二批由标定和方程求出：heavy\_from\_sample\_mol = 0.00005416666666667 mol；heavy\_from\_spike\_mol = 0.001800000000000 mol；total\_molecules\_mol = 0.007416666666667 mol。
\item 其余浓度、校准或统计结果：total\_heavy\_mol = 0.001854166666667 mol；total\_light\_mol = 0.005562500000000 mol；heavy\_light\_ratio = 0.3333333333333 ratio；fraction\_check = 0.2500000000000 fraction。
\item 用计算结果回代原观测量或响应方程，检查物料和计数归一以及标准态一致。
\item 当仪器存在背景漂移、分辨率限制、基体效应和非理想标定时，需重新评估系统偏差。
\item GitHub QA校核全部数值字段，但公式见证尚属于同源实现，不能称已独立验证。
\item 题目学术可信度与跨辑去重仍须专家逐题审核。
\end{enumerate}
\noindent\textbf{答案摘要}\par
sample\_mol = 0.00541666667 mol；spike\_mol = 0.002 mol；sample\_spike\_ratio = 2.70833333 ratio；heavy\_from\_sample\_mol = 5.41666667e-05 mol；heavy\_from\_spike\_mol = 0.0018 mol；total\_molecules\_mol = 0.00741666667 mol；total\_heavy\_mol = 0.00185416667 mol；total\_light\_mol = 0.0055625 mol；heavy\_light\_ratio = 0.333333333 ratio；fraction\_check = 0.25 fraction\par
\noindent\textbf{模型前提：}
理想分析仪器与信号计量，不代表未经评定的实验准确度。\par
\noindent\textcolor{red!65!black}{[跨辑去重未完成 · 草稿]} \par
\medskip\hrule\medskip
\Needspace{7\baselineskip}
\subsubsection*{UC-ADV15-A227 · 毛细管电泳迁移时间的电泳迁移率与电渗流分离}
\noindent\textbf{题目}\par
某毛细管电泳仪器总毛细管长度Lt=0.6000 m，检测点到入口距离Ld=0.5000 m，电压V=20000 V。中性电渗流标记物到达时间tEOF=100.0 s，单价分析离子到达时间tA=80.00 s，另一阴离子标记到达时间tB=125.0 s。电场沿管均匀，检测速度v=Ld/t，μapp=v/E=Lt Ld/(Vt)，μep=μapp−μEOF。求EOF表观迁移率、A/B表观迁移率、A/B本征电泳迁移率、A相对EOF速度比、两离子迁移率差、EOF体积流量未知时为何仍可求迁移率，以及A/B相对EOF的净迁移方向。离子没有吸附和电荷状态互变。\par
\medskip\noindent\textbf{分步解答}
\begin{enumerate}
\item 实验测量模型和给定条件：某毛细管电泳仪器总毛细管长度Lt=0.6000 m，检测点到入口距离Ld=0.5000 m，电压V=20000 V。中性电渗流标记物到达时间tEOF=100.0 s，单价分析离子到达时间tA=80.00 s，另一阴离子标记到达时间tB=125.0 s。电场沿管均匀，检测速度v=Ld/t，μapp=v/E=Lt Ld/(Vt)，μep=μapp−μEOF。求EOF表观迁移率、A/B表观迁移率、A/B本征电泳迁移率、A相对EOF速度比、两离子迁移率差、EOF体积流量未知时为何仍可求迁移率，以及A/B相对EOF的净迁移方向。离子没有吸附和电荷状态互变。
\item 完整反演或计量方程：μapp=LtLd/(V t)，μEOF=LtLd/(V tEOF)，μep=μapp−μEOF。
\item 确认原样、标准品和空白的质量或体积关系，所有光谱响应和电位基线按题面规定处理。
\item 第一批直接求得：mobility\_EOF = 1.500000000000e-7 m2/(Vs)；mobility\_A\_app = 1.875000000000e-7 m2/(Vs)；mobility\_B\_app = 1.200000000000e-7 m2/(Vs)。
\item 第二批由标定和方程求出：mobility\_A\_ep = 3.750000000000e-8 m2/(Vs)；mobility\_B\_ep = -3.000000000000e-8 m2/(Vs)；velocity\_A\_m\_s = 0.006250000000000 m/s。
\item 其余浓度、校准或统计结果：velocity\_B\_m\_s = 0.004000000000000 m/s；velocity\_EOF\_m\_s = 0.005000000000000 m/s；ep\_mobility\_difference = 6.750000000000e-8 m2/(Vs)；A\_over\_EOF\_velocity = 1.250000000000 ratio。
\item 用计算结果回代原观测量或响应方程，检查物料和计数归一以及标准态一致。
\item 当仪器存在背景漂移、分辨率限制、基体效应和非理想标定时，需重新评估系统偏差。
\item GitHub QA校核全部数值字段，但公式见证尚属于同源实现，不能称已独立验证。
\item 题目学术可信度与跨辑去重仍须专家逐题审核。
\end{enumerate}
\noindent\textbf{答案摘要}\par
mobility\_EOF = 1.5e-07 m2/(Vs)；mobility\_A\_app = 1.875e-07 m2/(Vs)；mobility\_B\_app = 1.2e-07 m2/(Vs)；mobility\_A\_ep = 3.75e-08 m2/(Vs)；mobility\_B\_ep = -3e-08 m2/(Vs)；velocity\_A\_m\_s = 0.00625 m/s；velocity\_B\_m\_s = 0.004 m/s；velocity\_EOF\_m\_s = 0.005 m/s；ep\_mobility\_difference = 6.75e-08 m2/(Vs)；A\_over\_EOF\_velocity = 1.25 ratio\par
\noindent\textbf{模型前提：}
理想分析仪器与信号计量，不代表未经评定的实验准确度。\par
\noindent\textcolor{red!65!black}{[跨辑去重未完成 · 草稿]} \par
\medskip\hrule\medskip
\Needspace{7\baselineskip}
\subsubsection*{UC-ADV15-A228 · 等温滴定量热法1:1结合焓的总热量与稀释空白}
\noindent\textbf{题目}\par
假想等温滴定量热实验中蛋白结合配体反应P+L⇌PL，某一完整滴定阶段按平衡模型已知新增结合量Δn=4.000×10⁻⁸ mol，反应结合焓ΔH=−30.00 kJ/mol，除结合外有定量稀释背景热Qblank=+0.1000 mJ，仪器测得积分热Qobs=Qbind+Qblank。另一个滴定阶段新增结合量为2.000×10⁻⁸ mol，具有相同焓和相同空白。求两阶段结合热、总观测热、扣空白热、两阶段观测热比、按扣空白反演的ΔH、把空白漏扣导致的表观ΔH、累计结合量与累计结合热。所有热统一为mJ，忽略热容和注射体积变化。\par
\medskip\noindent\textbf{分步解答}
\begin{enumerate}
\item 实验测量模型和给定条件：假想等温滴定量热实验中蛋白结合配体反应P+L⇌PL，某一完整滴定阶段按平衡模型已知新增结合量Δn=4.000×10⁻⁸ mol，反应结合焓ΔH=−30.00 kJ/mol，除结合外有定量稀释背景热Qblank=+0.1000 mJ，仪器测得积分热Qobs=Qbind+Qblank。另一个滴定阶段新增结合量为2.000×10⁻⁸ mol，具有相同焓和相同空白。求两阶段结合热、总观测热、扣空白热、两阶段观测热比、按扣空白反演的ΔH、把空白漏扣导致的表观ΔH、累计结合量与累计结合热。所有热统一为mJ，忽略热容和注射体积变化。
\item 完整反演或计量方程：Qbind=Δn ΔH；Qobs=Qbind+Qblank，ΔHapp=Qobs/Δn。
\item 确认原样、标准品和空白的质量或体积关系，所有光谱响应和电位基线按题面规定处理。
\item 第一批直接求得：binding\_heat1\_mJ = -1.200000000000 mJ；observed\_heat1\_mJ = -1.100000000000 mJ；binding\_heat2\_mJ = -0.6000000000000 mJ。
\item 第二批由标定和方程求出：observed\_heat2\_mJ = -0.5000000000000 mJ；blank\_mJ = 0.1000000000000 mJ；recovered\_binding\_enthalpy\_kJmol = -30.00000000000 kJ/mol。
\item 其余浓度、校准或统计结果：wrong\_enthalpy\_without\_blank\_kJmol = -27.50000000000 kJ/mol；observed\_heat\_ratio = 2.200000000000 ratio；cumulative\_bound\_mol = 6.000000000000e-8 mol；cumulative\_binding\_heat\_mJ = -1.800000000000 mJ。
\item 用计算结果回代原观测量或响应方程，检查物料和计数归一以及标准态一致。
\item 当仪器存在背景漂移、分辨率限制、基体效应和非理想标定时，需重新评估系统偏差。
\item GitHub QA校核全部数值字段，但公式见证尚属于同源实现，不能称已独立验证。
\item 题目学术可信度与跨辑去重仍须专家逐题审核。
\end{enumerate}
\noindent\textbf{答案摘要}\par
binding\_heat1\_mJ = -1.2 mJ；observed\_heat1\_mJ = -1.1 mJ；binding\_heat2\_mJ = -0.6 mJ；observed\_heat2\_mJ = -0.5 mJ；blank\_mJ = 0.1 mJ；recovered\_binding\_enthalpy\_kJmol = -30 kJ/mol；wrong\_enthalpy\_without\_blank\_kJmol = -27.5 kJ/mol；observed\_heat\_ratio = 2.2 ratio；cumulative\_bound\_mol = 6e-08 mol；cumulative\_binding\_heat\_mJ = -1.8 mJ\par
\noindent\textbf{模型前提：}
理想分析仪器与信号计量，不代表未经评定的实验准确度。\par
\noindent\textcolor{red!65!black}{[跨辑去重未完成 · 草稿]} \par
\medskip\hrule\medskip
\Needspace{7\baselineskip}
\subsubsection*{UC-ADV15-A229 · 质谱峰自然丰度中的同位素卷积二项分布}
\noindent\textbf{题目}\par
某分子含两个相互独立等价氯原子，每个氯原子有35Cl天然丰度p=0.7578、37Cl丰度q=0.2422，其他元素的同位素丰度在本题忽略。两氯分子呈M、M+2、M+4三个质谱同位素峰，几率分别为p²、2pq、q²。假设响应效率不依赖同位素，求三峰概率、相对于M峰的强度比、其相对M=100基准的峰强度、观察100000个分子时各峰预期计数、相对分布均值所对应的37Cl个数及每分子37Cl个数方差。说明中间峰的组合因子2来自两个可替换氯位点。\par
\medskip\noindent\textbf{分步解答}
\begin{enumerate}
\item 实验测量模型和给定条件：某分子含两个相互独立等价氯原子，每个氯原子有35Cl天然丰度p=0.7578、37Cl丰度q=0.2422，其他元素的同位素丰度在本题忽略。两氯分子呈M、M+2、M+4三个质谱同位素峰，几率分别为p²、2pq、q²。假设响应效率不依赖同位素，求三峰概率、相对于M峰的强度比、其相对M=100基准的峰强度、观察100000个分子时各峰预期计数、相对分布均值所对应的37Cl个数及每分子37Cl个数方差。说明中间峰的组合因子2来自两个可替换氯位点。
\item 完整反演或计量方程：PM=p²，PM2=2pq，PM4=q²；E[N37]=2q，Var[N37]=2pq。
\item 确认原样、标准品和空白的质量或体积关系，所有光谱响应和电位基线按题面规定处理。
\item 第一批直接求得：prob\_M = 0.5742608400000 fraction；prob\_Mplus2 = 0.3670783200000 fraction；prob\_Mplus4 = 0.05866084000000 fraction。
\item 第二批由标定和方程求出：ratio\_M2\_M = 0.6392187912378 ratio；ratio\_M4\_M = 0.1021501657679 ratio；relative\_Mplus2\_intensity\_to100 = 63.92187912378 counts\_relative。
\item 其余浓度、校准或统计结果：relative\_Mplus4\_intensity\_to100 = 10.21501657679 counts\_relative；expected\_counts\_Mplus2 = 36707.83200000 counts；mean\_37Cl\_atoms = 0.4844000000000 atoms；variance\_37Cl\_atoms = 0.3670783200000 atoms2。
\item 用计算结果回代原观测量或响应方程，检查物料和计数归一以及标准态一致。
\item 当仪器存在背景漂移、分辨率限制、基体效应和非理想标定时，需重新评估系统偏差。
\item GitHub QA校核全部数值字段，但公式见证尚属于同源实现，不能称已独立验证。
\item 题目学术可信度与跨辑去重仍须专家逐题审核。
\end{enumerate}
\noindent\textbf{答案摘要}\par
prob\_M = 0.57426084 fraction；prob\_Mplus2 = 0.36707832 fraction；prob\_Mplus4 = 0.05866084 fraction；ratio\_M2\_M = 0.639218791 ratio；ratio\_M4\_M = 0.102150166 ratio；relative\_Mplus2\_intensity\_to100 = 63.9218791 counts\_relative；relative\_Mplus4\_intensity\_to100 = 10.2150166 counts\_relative；expected\_counts\_Mplus2 = 36707.832 counts；mean\_37Cl\_atoms = 0.4844 atoms；variance\_37Cl\_atoms = 0.36707832 atoms2\par
\noindent\textbf{模型前提：}
理想分析仪器与信号计量，不代表未经评定的实验准确度。\par
\noindent\textcolor{red!65!black}{[跨辑去重未完成 · 草稿]} \par
\medskip\hrule\medskip
\Needspace{7\baselineskip}
\subsubsection*{UC-ADV15-A230 · 定量电化学传感器Michaelis–Menten饱和电流的Kₘ反演}
\noindent\textbf{题目}\par
某固定化酶电化学传感器的稳态催化电流i=i\_max C/(K\_m+C)，其中分析物浓度C用mmol/L计。给定在C1=1.000 mmol/L测得i1=2.000 μA，在C2=3.000 mmol/L测得i2=4.000 μA，假设无背景、非特异吸附与传质限制，电极电流完全由一级酶催化机制控制。通过Hanes形式C/i=(K\_m+C)/i\_max或双点直接求Km及imax，预测C=5.000、10.00时的电流、初始低浓度灵敏度imax/Km、半饱和浓度、两个电流比和最大电流归一化分率。注意高浓度不能继续按低浓度线性标定。\par
\medskip\noindent\textbf{分步解答}
\begin{enumerate}
\item 实验测量模型和给定条件：某固定化酶电化学传感器的稳态催化电流i=i\_max C/(K\_m+C)，其中分析物浓度C用mmol/L计。给定在C1=1.000 mmol/L测得i1=2.000 μA，在C2=3.000 mmol/L测得i2=4.000 μA，假设无背景、非特异吸附与传质限制，电极电流完全由一级酶催化机制控制。通过Hanes形式C/i=(K\_m+C)/i\_max或双点直接求Km及imax，预测C=5.000、10.00时的电流、初始低浓度灵敏度imax/Km、半饱和浓度、两个电流比和最大电流归一化分率。注意高浓度不能继续按低浓度线性标定。
\item 完整反演或计量方程：1/i=(K\_m/i\_max)(1/C)+1/i\_max；两点求得K\_m=3.000 mM、i\_max=8.000 μA。
\item 确认原样、标准品和空白的质量或体积关系，所有光谱响应和电位基线按题面规定处理。
\item 第一批直接求得：Km\_mM = 3.000000000000 mmol/L；imax\_microA = 8.000000000000 microA；current\_C1\_microA = 2.000000000000 microA。
\item 第二批由标定和方程求出：current\_C3\_microA = 4.000000000000 microA；current\_C5\_microA = 5.000000000000 microA；current\_C10\_microA = 6.153846153846 microA。
\item 其余浓度、校准或统计结果：initial\_sensitivity\_microA\_mM = 2.666666666667 microA/mM；half\_saturation\_C\_mM = 3.000000000000 mmol/L；ratio\_i10\_over\_i5 = 1.230769230769 ratio；fraction\_max\_at\_C10 = 0.7692307692308 fraction。
\item 用计算结果回代原观测量或响应方程，检查物料和计数归一以及标准态一致。
\item 当仪器存在背景漂移、分辨率限制、基体效应和非理想标定时，需重新评估系统偏差。
\item GitHub QA校核全部数值字段，但公式见证尚属于同源实现，不能称已独立验证。
\item 题目学术可信度与跨辑去重仍须专家逐题审核。
\end{enumerate}
\noindent\textbf{答案摘要}\par
Km\_mM = 3 mmol/L；imax\_microA = 8 microA；current\_C1\_microA = 2 microA；current\_C3\_microA = 4 microA；current\_C5\_microA = 5 microA；current\_C10\_microA = 6.15384615 microA；initial\_sensitivity\_microA\_mM = 2.66666667 microA/mM；half\_saturation\_C\_mM = 3 mmol/L；ratio\_i10\_over\_i5 = 1.23076923 ratio；fraction\_max\_at\_C10 = 0.769230769 fraction\par
\noindent\textbf{模型前提：}
理想分析仪器与信号计量，不代表未经评定的实验准确度。\par
\noindent\textcolor{red!65!black}{[跨辑去重未完成 · 草稿]} \par
\medskip\hrule\medskip
\Needspace{7\baselineskip}
\subsubsection*{UC-ADV15-A231 · 等离子体原子吸收光谱双线校正的背景吸光反演}
\noindent\textbf{题目}\par
某原子吸收光谱AAS在分析线λa和邻近非吸收参比线λr同时记录样品吸光度，理想模型Aa=kC+B、Ar=B，其中k=0.5000 L/mg，B是宽带背景吸光，未知样本Aa=0.8500、Ar=0.1000。加入1:1稀释、并假定背景吸光随样本浓度同倍减半而新信号服从同一模型。求样本扣背景吸光、原样分析物浓度(mg/L)、原样背景、稀释样Aa与Ar、稀释后扣背景吸光及浓度、忽略背景时的浓度与相对高估，以及原样信号由分析物贡献的比例。说明参比线不能出现分析物吸收，否则两线扣除有偏。\par
\medskip\noindent\textbf{分步解答}
\begin{enumerate}
\item 实验测量模型和给定条件：某原子吸收光谱AAS在分析线λa和邻近非吸收参比线λr同时记录样品吸光度，理想模型Aa=kC+B、Ar=B，其中k=0.5000 L/mg，B是宽带背景吸光，未知样本Aa=0.8500、Ar=0.1000。加入1:1稀释、并假定背景吸光随样本浓度同倍减半而新信号服从同一模型。求样本扣背景吸光、原样分析物浓度(mg/L)、原样背景、稀释样Aa与Ar、稀释后扣背景吸光及浓度、忽略背景时的浓度与相对高估，以及原样信号由分析物贡献的比例。说明参比线不能出现分析物吸收，否则两线扣除有偏。
\item 完整反演或计量方程：C=(Aa−Ar)/k；稀释后C'=C/2、B'=B/2。
\item 确认原样、标准品和空白的质量或体积关系，所有光谱响应和电位基线按题面规定处理。
\item 第一批直接求得：analyte\_absorbance = 0.7500000000000 AU；original\_concentration\_mgL = 1.500000000000 mg/L；background\_absorbance = 0.1000000000000 AU。
\item 第二批由标定和方程求出：diluted\_analytical\_absorbance = 0.4250000000000 AU；diluted\_reference\_absorbance = 0.05000000000000 AU；diluted\_corrected\_absorbance = 0.3750000000000 AU。
\item 其余浓度、校准或统计结果：diluted\_concentration\_mgL = 0.7500000000000 mg/L；naive\_uncorrected\_concentration\_mgL = 1.700000000000 mg/L；naive\_relative\_overestimate = 0.1333333333333 fraction；analyte\_signal\_fraction = 0.8823529411765 fraction。
\item 用计算结果回代原观测量或响应方程，检查物料和计数归一以及标准态一致。
\item 当仪器存在背景漂移、分辨率限制、基体效应和非理想标定时，需重新评估系统偏差。
\item GitHub QA校核全部数值字段，但公式见证尚属于同源实现，不能称已独立验证。
\item 题目学术可信度与跨辑去重仍须专家逐题审核。
\end{enumerate}
\noindent\textbf{答案摘要}\par
analyte\_absorbance = 0.75 AU；original\_concentration\_mgL = 1.5 mg/L；background\_absorbance = 0.1 AU；diluted\_analytical\_absorbance = 0.425 AU；diluted\_reference\_absorbance = 0.05 AU；diluted\_corrected\_absorbance = 0.375 AU；diluted\_concentration\_mgL = 0.75 mg/L；naive\_uncorrected\_concentration\_mgL = 1.7 mg/L；naive\_relative\_overestimate = 0.133333333 fraction；analyte\_signal\_fraction = 0.882352941 fraction\par
\noindent\textbf{模型前提：}
理想分析仪器与信号计量，不代表未经评定的实验准确度。\par
\noindent\textcolor{red!65!black}{[跨辑去重未完成 · 草稿]} \par
\medskip\hrule\medskip
\Needspace{7\baselineskip}
\subsubsection*{UC-ADV15-A232 · 气相傅里叶红外光谱线强度温度校正与浓度误差}
\noindent\textbf{题目}\par
某气相FTIR吸收测量模型A(T)=S(T)cL，光程L=0.1000 m、浓度c以任意固定mol/m³计。单吸收谱线强度温度依赖由题设近似S(T)=S300 (300/T)exp[−E''/R(1/T−1/300)]，其中S300=0.5000 m²/mol，低能级能E''=5000 J/mol。样品在T=350.0 K实际浓度c=0.02000 mol/m³。计算S350、实际吸光度A350、若错误用S300标定得到的表观c、相对系统偏差、温度修正系数S350/S300、若实际光程加倍后A值以及300K在相同浓度光程下A300；再求两温度吸光度比与反演原始浓度的结果。模型不含压力展宽、谱线重叠和分子配分函数更复杂项。\par
\medskip\noindent\textbf{分步解答}
\begin{enumerate}
\item 实验测量模型和给定条件：某气相FTIR吸收测量模型A(T)=S(T)cL，光程L=0.1000 m、浓度c以任意固定mol/m³计。单吸收谱线强度温度依赖由题设近似S(T)=S300 (300/T)exp[−E''/R(1/T−1/300)]，其中S300=0.5000 m²/mol，低能级能E''=5000 J/mol。样品在T=350.0 K实际浓度c=0.02000 mol/m³。计算S350、实际吸光度A350、若错误用S300标定得到的表观c、相对系统偏差、温度修正系数S350/S300、若实际光程加倍后A值以及300K在相同浓度光程下A300；再求两温度吸光度比与反演原始浓度的结果。模型不含压力展宽、谱线重叠和分子配分函数更复杂项。
\item 完整反演或计量方程：S350=S300(300/350)exp[−(5000/R)(1/350−1/300)]；capp=A350/(S300 L)。
\item 确认原样、标准品和空白的质量或体积关系，所有光谱响应和电位基线按题面规定处理。
\item 第一批直接求得：line\_strength\_S350 = 0.5706751700448 m2/mol；absorbance\_350 = 0.001141350340090 AU；apparent\_c\_using300 = 0.02282700680179 mol/m3。
\item 第二批由标定和方程求出：relative\_bias = 0.1413503400895 fraction；line\_strength\_ratio = 1.141350340090 ratio；absorbance\_double\_path = 0.002282700680179 AU。
\item 其余浓度、校准或统计结果：absorbance\_300 = 0.001000000000000 AU；absorbance\_ratio350\_300 = 1.141350340090 ratio；recovered\_actual\_c = 0.02000000000000 mol/m3；path\_length\_m = 0.1000000000000 m。
\item 用计算结果回代原观测量或响应方程，检查物料和计数归一以及标准态一致。
\item 当仪器存在背景漂移、分辨率限制、基体效应和非理想标定时，需重新评估系统偏差。
\item GitHub QA校核全部数值字段，但公式见证尚属于同源实现，不能称已独立验证。
\item 题目学术可信度与跨辑去重仍须专家逐题审核。
\end{enumerate}
\noindent\textbf{答案摘要}\par
line\_strength\_S350 = 0.57067517 m2/mol；absorbance\_350 = 0.00114135034 AU；apparent\_c\_using300 = 0.0228270068 mol/m3；relative\_bias = 0.14135034 fraction；line\_strength\_ratio = 1.14135034 ratio；absorbance\_double\_path = 0.00228270068 AU；absorbance\_300 = 0.001 AU；absorbance\_ratio350\_300 = 1.14135034 ratio；recovered\_actual\_c = 0.02 mol/m3；path\_length\_m = 0.1 m\par
\noindent\textbf{模型前提：}
理想分析仪器与信号计量，不代表未经评定的实验准确度。\par
\noindent\textcolor{red!65!black}{[跨辑去重未完成 · 草稿]} \par
\medskip\hrule\medskip
\Needspace{7\baselineskip}
\subsubsection*{UC-ADV15-A233 · 电化学阻抗Randles串联溶液电阻与并联RC半圆}
\noindent\textbf{题目}\par
某电极等效Randles电路由溶液电阻Rs=10.00 Ω与电荷转移电阻Rct=100.0 Ω并联双电层电容Cdl=100.0 μF串联组成，假设无Warburg扩散。阻抗Z(ω)=Rs+Rct/(1+iωRct Cdl)，以Z=Z'−iZ''（Z''为正损耗）表示。分别在f=1.000 Hz和f=100.0 Hz求实部Z'和正损耗Z''、半圆顶点频率fmax=1/(2πRct Cdl)、半圆两端高低频电阻、顶点阻抗实部及损耗峰值、100 Hz阻抗模和低频/高频电阻比。数值全部用Ω，电容μF须转换为F。\par
\medskip\noindent\textbf{分步解答}
\begin{enumerate}
\item 实验测量模型和给定条件：某电极等效Randles电路由溶液电阻Rs=10.00 Ω与电荷转移电阻Rct=100.0 Ω并联双电层电容Cdl=100.0 μF串联组成，假设无Warburg扩散。阻抗Z(ω)=Rs+Rct/(1+iωRct Cdl)，以Z=Z'−iZ''（Z''为正损耗）表示。分别在f=1.000 Hz和f=100.0 Hz求实部Z'和正损耗Z''、半圆顶点频率fmax=1/(2πRct Cdl)、半圆两端高低频电阻、顶点阻抗实部及损耗峰值、100 Hz阻抗模和低频/高频电阻比。数值全部用Ω，电容μF须转换为F。
\item 完整反演或计量方程：Z'=Rs+Rct/(1+(ωτ)²)，Z''=Rctωτ/(1+(ωτ)²)，τ=RctCdl。
\item 确认原样、标准品和空白的质量或体积关系，所有光谱响应和电位基线按题面规定处理。
\item 第一批直接求得：real\_Z\_1Hz = 109.6067682407 ohm；loss\_Z\_1Hz = 6.258477827057 ohm；real\_Z\_100Hz = 12.47045230319 ohm。
\item 第二批由标定和方程求出：loss\_Z\_100Hz = 15.52230961346 ohm；semicircle\_peak\_frequency\_Hz = 15.91549430919 Hz；low\_frequency\_resistance = 110.0000000000 ohm。
\item 其余浓度、校准或统计结果：high\_frequency\_resistance = 10.00000000000 ohm；real\_Z\_at\_peak = 60.00000000000 ohm；loss\_Z\_at\_peak = 50.00000000000 ohm；impedance\_magnitude\_100Hz = 19.91115959411 ohm。
\item 用计算结果回代原观测量或响应方程，检查物料和计数归一以及标准态一致。
\item 当仪器存在背景漂移、分辨率限制、基体效应和非理想标定时，需重新评估系统偏差。
\item GitHub QA校核全部数值字段，但公式见证尚属于同源实现，不能称已独立验证。
\item 题目学术可信度与跨辑去重仍须专家逐题审核。
\end{enumerate}
\noindent\textbf{答案摘要}\par
real\_Z\_1Hz = 109.606768 ohm；loss\_Z\_1Hz = 6.25847783 ohm；real\_Z\_100Hz = 12.4704523 ohm；loss\_Z\_100Hz = 15.5223096 ohm；semicircle\_peak\_frequency\_Hz = 15.9154943 Hz；low\_frequency\_resistance = 110 ohm；high\_frequency\_resistance = 10 ohm；real\_Z\_at\_peak = 60 ohm；loss\_Z\_at\_peak = 50 ohm；impedance\_magnitude\_100Hz = 19.9111596 ohm\par
\noindent\textbf{模型前提：}
理想分析仪器与信号计量，不代表未经评定的实验准确度。\par
\noindent\textcolor{red!65!black}{[跨辑去重未完成 · 草稿]} \par
\medskip\hrule\medskip
\Needspace{7\baselineskip}
\subsubsection*{UC-ADV15-A234 · 离子选择电极双点校准的膜电势与未知活度}
\noindent\textbf{题目}\par
298.15 K某选择性电极测一价阳离子X⁺，理想Nernst形式E=E°+(RT/F)ln aX，E°=0.3000 V。标准液活动度a1=0.001000、a2=0.01000（以1 M归一化），未知样的真实活动度aU=0.003000，所有活度系数与参比电位已经标定、无干扰离子。求两标准的电位、未知电位、由电位反算aU、未知pX、每十倍浓度电位斜率、未知点在两标准之间的对数插值分率、未知aX加倍时电极电位和电位增量、两标准电压差。注意绝对电位与参比标定有关，但差分斜率与Nernst一致。\par
\medskip\noindent\textbf{分步解答}
\begin{enumerate}
\item 实验测量模型和给定条件：298.15 K某选择性电极测一价阳离子X⁺，理想Nernst形式E=E°+(RT/F)ln aX，E°=0.3000 V。标准液活动度a1=0.001000、a2=0.01000（以1 M归一化），未知样的真实活动度aU=0.003000，所有活度系数与参比电位已经标定、无干扰离子。求两标准的电位、未知电位、由电位反算aU、未知pX、每十倍浓度电位斜率、未知点在两标准之间的对数插值分率、未知aX加倍时电极电位和电位增量、两标准电压差。注意绝对电位与参比标定有关，但差分斜率与Nernst一致。
\item 完整反演或计量方程：E=E°+(RT/F)ln a，a=exp[(E−E°)F/RT]。
\item 确认原样、标准品和空白的质量或体积关系，所有光谱响应和电位基线按题面规定处理。
\item 第一批直接求得：E\_standard1 = 0.1225219509428 V；E\_standard2 = 0.1816813006286 V；E\_unknown = 0.1507481340933 V。
\item 第二批由标定和方程求出：activity\_recovered = 0.003000000000000 activity；pX = 2.522878745280 pX；slope\_per\_decade = 0.05915934968572 V。
\item 其余浓度、校准或统计结果：log\_interpolation\_fraction = 0.4771212547197 fraction；E\_double\_activity = 0.1685568728727 V；potential\_increment\_double = 0.01780873877938 V；standard\_voltage\_difference = 0.05915934968572 V。
\item 用计算结果回代原观测量或响应方程，检查物料和计数归一以及标准态一致。
\item 当仪器存在背景漂移、分辨率限制、基体效应和非理想标定时，需重新评估系统偏差。
\item GitHub QA校核全部数值字段，但公式见证尚属于同源实现，不能称已独立验证。
\item 题目学术可信度与跨辑去重仍须专家逐题审核。
\end{enumerate}
\noindent\textbf{答案摘要}\par
E\_standard1 = 0.122521951 V；E\_standard2 = 0.181681301 V；E\_unknown = 0.150748134 V；activity\_recovered = 0.003 activity；pX = 2.52287875 pX；slope\_per\_decade = 0.0591593497 V；log\_interpolation\_fraction = 0.477121255 fraction；E\_double\_activity = 0.168556873 V；potential\_increment\_double = 0.0178087388 V；standard\_voltage\_difference = 0.0591593497 V\par
\noindent\textbf{模型前提：}
理想分析仪器与信号计量，不代表未经评定的实验准确度。\par
\noindent\textcolor{red!65!black}{[跨辑去重未完成 · 草稿]} \par
\medskip\hrule\medskip
\subsection*{有机化学}
\addcontentsline{toc}{subsection}{第15辑 / 有机化学}
\Needspace{7\baselineskip}
\subsubsection*{UC-ADV15-O223 · 手性醇酶促动力学拆分的转化率与对映体过量}
\noindent\textbf{题目}\par
某外消旋二级醇由R、S对映异构体各0.005000 mol/L构成，分别按不可逆一级反应转化为对应酯PR、PS，反应速度kR=0.04000 s⁻¹、kS=0.01000 s⁻¹，反应过程中不外消旋、无产物水解或逆反应。在t=30.00 s计算残余R、S浓度和各产物浓度、总底物转化率、剩余醇的eeS=(S−R)/(S+R)、酯产物eeR=(PR−PS)/(PR+PS)，以及残余醇/产物质量守恒、瞬时初始选择性比和产物总浓度。请说明两种ee与总体转化率之间并非独立，且不应把残余底物的ee视为产物ee。\par
\medskip\noindent\textbf{分步解答}
\begin{enumerate}
\item 完整有机化学反应与初始组成：某外消旋二级醇由R、S对映异构体各0.005000 mol/L构成，分别按不可逆一级反应转化为对应酯PR、PS，反应速度kR=0.04000 s⁻¹、kS=0.01000 s⁻¹，反应过程中不外消旋、无产物水解或逆反应。在t=30.00 s计算残余R、S浓度和各产物浓度、总底物转化率、剩余醇的eeS=(S−R)/(S+R)、酯产物eeR=(PR−PS)/(PR+PS)，以及残余醇/产物质量守恒、瞬时初始选择性比和产物总浓度。请说明两种ee与总体转化率之间并非独立，且不应把残余底物的ee视为产物ee。
\item 从反应物库存、速率网络或统计假设建立数学式：R=R0e\textasciicircum{}(−kR t)，S=S0e\textasciicircum{}(−kS t)，PR=R0−R，PS=S0−S；X=(PR+PS)/(R0+S0)。
\item 写出底物与产物的严格物质守恒或反应构型概率，注意一级和二级速率的量纲转换。
\item 根据模型计算初始常数和分率：remaining\_R\_M = 0.0015059710595610 mol/L；remaining\_S\_M = 0.0037040911034086 mol/L；product\_R\_M = 0.0034940289404390 mol/L。
\item 继续计算反应进行后的浓度或分布：product\_S\_M = 0.0012959088965914 mol/L；conversion = 0.47899378370304 fraction；substrate\_ee\_S = 0.42189900525001 fraction。
\item 得到其余转化、选择性或分散度量：product\_ee\_R = 0.45890366819673 fraction；total\_product\_M = 0.0047899378370304 mol/L；molecule\_inventory\_M = 0.010000000000000 mol/L；initial\_selectivity\_kR\_kS = 4.0000000000000 ratio。
\item 代回底物+产物的骨架守恒、概率归一或Eyring速率比，确认所有答案范围和符号。
\item 在长时间极限、没有反应极限及温度变化时，检查公式预言的物理趋势。
\item 本题所忽略的逆反应、构象重排、聚合动力学记忆和溶剂效应不应从理想数值随意外推。
\item 数值见证脚本验证输出一致性，独立实现和跨辑模型原创性仍需人工复核。
\end{enumerate}
\noindent\textbf{答案摘要}\par
remaining\_R\_M = 0.00150597106 mol/L；remaining\_S\_M = 0.0037040911 mol/L；product\_R\_M = 0.00349402894 mol/L；product\_S\_M = 0.0012959089 mol/L；conversion = 0.478993784 fraction；substrate\_ee\_S = 0.421899005 fraction；product\_ee\_R = 0.458903668 fraction；total\_product\_M = 0.00478993784 mol/L；molecule\_inventory\_M = 0.01 mol/L；initial\_selectivity\_kR\_kS = 4 ratio\par
\noindent\textbf{模型前提：}
本题无副反应和理想配分条件为假设；需独立科学审查。\par
\noindent\textcolor{red!65!black}{[跨辑去重未完成 · 草稿]} \par
\medskip\hrule\medskip
\Needspace{7\baselineskip}
\subsubsection*{UC-ADV15-O224 · 有机羰基化Curtin–Hammett构象通量与产物选择性}
\noindent\textbf{题目}\par
某反应底物A/B两构象快速互变，在298.15 K热平衡下[A]:[B]=1:3。两构象分别经不可逆一级过渡态生成不同产物PA、PB，且各自反应常数kA=0.2000 s⁻¹、kB=0.04000 s⁻¹，构象互变始终远快于化学消耗，总未反应底物初始0.02000 M。计算构象A/B分率、两途径有效一级速率、总观察速率、两产物分支比和PA/PB比、反应t=20.00 s后残余底物与两产物浓度及半衰期。解释即使A为少数构象也可贡献主要产物，不能单凭构象基态分率判断产物组成。\par
\medskip\noindent\textbf{分步解答}
\begin{enumerate}
\item 完整有机化学反应与初始组成：某反应底物A/B两构象快速互变，在298.15 K热平衡下[A]:[B]=1:3。两构象分别经不可逆一级过渡态生成不同产物PA、PB，且各自反应常数kA=0.2000 s⁻¹、kB=0.04000 s⁻¹，构象互变始终远快于化学消耗，总未反应底物初始0.02000 M。计算构象A/B分率、两途径有效一级速率、总观察速率、两产物分支比和PA/PB比、反应t=20.00 s后残余底物与两产物浓度及半衰期。解释即使A为少数构象也可贡献主要产物，不能单凭构象基态分率判断产物组成。
\item 从反应物库存、速率网络或统计假设建立数学式：kobs=fA kA+fB kB；bA=fA kA/kobs；PA=C0(1−exp(−kobs t))bA。
\item 写出底物与产物的严格物质守恒或反应构型概率，注意一级和二级速率的量纲转换。
\item 根据模型计算初始常数和分率：fraction\_A = 0.25000000000000 fraction；fraction\_B = 0.75000000000000 fraction；flux\_rate\_A = 0.050000000000000 s\textasciicircum{}-1。
\item 继续计算反应进行后的浓度或分布：flux\_rate\_B = 0.030000000000000 s\textasciicircum{}-1；kobs = 0.080000000000000 s\textasciicircum{}-1；branch\_product\_A = 0.62500000000000 fraction。
\item 得到其余转化、选择性或分散度量：branch\_product\_B = 0.37500000000000 fraction；product\_A\_B\_ratio = 1.6666666666667 ratio；remaining\_substrate\_20 = 0.0040379303598931 mol/L；product\_A\_20 = 0.0099762935250668 mol/L；product\_B\_20 = 0.0059857761150401 mol/L；half\_life = 8.6643397569993 s。
\item 代回底物+产物的骨架守恒、概率归一或Eyring速率比，确认所有答案范围和符号。
\item 在长时间极限、没有反应极限及温度变化时，检查公式预言的物理趋势。
\item 本题所忽略的逆反应、构象重排、聚合动力学记忆和溶剂效应不应从理想数值随意外推。
\item 数值见证脚本验证输出一致性，独立实现和跨辑模型原创性仍需人工复核。
\end{enumerate}
\noindent\textbf{答案摘要}\par
fraction\_A = 0.25 fraction；fraction\_B = 0.75 fraction；flux\_rate\_A = 0.05 s\textasciicircum{}-1；flux\_rate\_B = 0.03 s\textasciicircum{}-1；kobs = 0.08 s\textasciicircum{}-1；branch\_product\_A = 0.625 fraction；branch\_product\_B = 0.375 fraction；product\_A\_B\_ratio = 1.66666667 ratio；remaining\_substrate\_20 = 0.00403793036 mol/L；product\_A\_20 = 0.00997629353 mol/L；product\_B\_20 = 0.00598577612 mol/L；half\_life = 8.66433976 s\par
\noindent\textbf{模型前提：}
本题无副反应和理想配分条件为假设；需独立科学审查。\par
\noindent\textcolor{red!65!black}{[跨辑去重未完成 · 草稿]} \par
\medskip\hrule\medskip
\Needspace{7\baselineskip}
\subsubsection*{UC-ADV15-O225 · 烯烃催化串联异构化–加氢的竞争中间体峰值}
\noindent\textbf{题目}\par
某有机烯烃A先以一级k1=0.08000 s⁻¹异构化为B，再以一级k2=0.02000 s⁻¹不可逆加氢生成饱和产物C，无其它通道，初始[A]=0.01000 M。求t=20.00 s时三物种浓度、各物种分率、B达到最大值的时间和峰值浓度、峰值时A与B两通道的瞬时反应速率以及最终全部骨架转化为C的平衡上限。必须从微分方程dA/dt=−k1A、dB/dt=k1A−k2B求解，不能简单用两个独立指数下降近似B。\par
\medskip\noindent\textbf{分步解答}
\begin{enumerate}
\item 完整有机化学反应与初始组成：某有机烯烃A先以一级k1=0.08000 s⁻¹异构化为B，再以一级k2=0.02000 s⁻¹不可逆加氢生成饱和产物C，无其它通道，初始[A]=0.01000 M。求t=20.00 s时三物种浓度、各物种分率、B达到最大值的时间和峰值浓度、峰值时A与B两通道的瞬时反应速率以及最终全部骨架转化为C的平衡上限。必须从微分方程dA/dt=−k1A、dB/dt=k1A−k2B求解，不能简单用两个独立指数下降近似B。
\item 从反应物库存、速率网络或统计假设建立数学式：A=C0e\textasciicircum{}(−k1t)，B=C0 k1/(k1−k2)[e\textasciicircum{}(−k2t)−e\textasciicircum{}(−k1t)]，C=C0−A−B。
\item 写出底物与产物的严格物质守恒或反应构型概率，注意一级和二级速率的量纲转换。
\item 根据模型计算初始常数和分率：A\_20\_M = 0.0020189651799466 mol/L；B\_20\_M = 0.0062456470405465 mol/L；C\_20\_M = 0.0017353877795070 mol/L。
\item 继续计算反应进行后的浓度或分布：fraction\_A = 0.20189651799466 fraction；fraction\_B = 0.62456470405465 fraction；fraction\_C = 0.17353877795070 fraction。
\item 得到其余转化、选择性或分散度量：B\_peak\_time\_s = 23.104906018665 s；B\_peak\_concentration\_M = 0.0062996052494744 mol/L；formation\_rate\_at\_B\_peak = 0.00012599210498949 mol/(Ls)；consumption\_rate\_at\_B\_peak = 0.00012599210498949 mol/(Ls)；total\_inventory\_M = 0.010000000000000 mol/L。
\item 代回底物+产物的骨架守恒、概率归一或Eyring速率比，确认所有答案范围和符号。
\item 在长时间极限、没有反应极限及温度变化时，检查公式预言的物理趋势。
\item 本题所忽略的逆反应、构象重排、聚合动力学记忆和溶剂效应不应从理想数值随意外推。
\item 数值见证脚本验证输出一致性，独立实现和跨辑模型原创性仍需人工复核。
\end{enumerate}
\noindent\textbf{答案摘要}\par
A\_20\_M = 0.00201896518 mol/L；B\_20\_M = 0.00624564704 mol/L；C\_20\_M = 0.00173538778 mol/L；fraction\_A = 0.201896518 fraction；fraction\_B = 0.624564704 fraction；fraction\_C = 0.173538778 fraction；B\_peak\_time\_s = 23.104906 s；B\_peak\_concentration\_M = 0.00629960525 mol/L；formation\_rate\_at\_B\_peak = 0.000125992105 mol/(Ls)；consumption\_rate\_at\_B\_peak = 0.000125992105 mol/(Ls)；total\_inventory\_M = 0.01 mol/L\par
\noindent\textbf{模型前提：}
本题无副反应和理想配分条件为假设；需独立科学审查。\par
\noindent\textcolor{red!65!black}{[跨辑去重未完成 · 草稿]} \par
\medskip\hrule\medskip
\Needspace{7\baselineskip}
\subsubsection*{UC-ADV15-O226 · 酯化竞争中极性溶剂促进离解的Hammond温度关联}
\noindent\textbf{题目}\par
某芳香酯发生两条不可逆平行裂解路径生成不同产物P1、P2，分别具有活化焓H‡1=55.00、H‡2=60.00 kJ/mol，活化熵S‡1=−40.00、S‡2=−25.00 J/(mol K)，两条途径各服从相同Eyring形式k\_i=kBT/h exp(S‡i/R−H‡i/RT)。温度为298.15与350.0 K。计算两个温度下两路径的速率常数、k1/k2分支比、P1产物分率、活化自由能差G1−G2、两个温度下产物比例变化与温度上升导致的路径1速率倍数。取kB=1.380649×10⁻²³ J/K、h=6.62607015×10⁻³⁴ J·s，忽略逆反应及熵焓随温度变化。\par
\medskip\noindent\textbf{分步解答}
\begin{enumerate}
\item 完整有机化学反应与初始组成：某芳香酯发生两条不可逆平行裂解路径生成不同产物P1、P2，分别具有活化焓H‡1=55.00、H‡2=60.00 kJ/mol，活化熵S‡1=−40.00、S‡2=−25.00 J/(mol K)，两条途径各服从相同Eyring形式k\_i=kBT/h exp(S‡i/R−H‡i/RT)。温度为298.15与350.0 K。计算两个温度下两路径的速率常数、k1/k2分支比、P1产物分率、活化自由能差G1−G2、两个温度下产物比例变化与温度上升导致的路径1速率倍数。取kB=1.380649×10⁻²³ J/K、h=6.62607015×10⁻³⁴ J·s，忽略逆反应及熵焓随温度变化。
\item 从反应物库存、速率网络或统计假设建立数学式：ki=(kBT/h) exp(Si/R−Hi/RT)，f1=k1/(k1+k2)，ΔG1−2=(H1−H2)−T(S1−S2)。
\item 写出底物与产物的严格物质守恒或反应构型概率，注意一级和二级速率的量纲转换。
\item 根据模型计算初始常数和分率：k1\_298 = 11.703995439443 s\textasciicircum{}-1；k2\_298 = 9.4596777920013 s\textasciicircum{}-1；k1\_350 = 367.62935978388 s\textasciicircum{}-1。
\item 继续计算反应进行后的浓度或分布：k2\_350 = 400.60827565140 s\textasciicircum{}-1；branch\_ratio\_298 = 1.2372509610570 ratio；branch\_ratio\_350 = 0.91767789665877 ratio。
\item 得到其余转化、选择性或分散度量：P1\_fraction\_298 = 0.55302287610705 fraction；P1\_fraction\_350 = 0.47853599306623 fraction；deltaG1\_minus2\_298\_kJ = -0.52775000000000 kJ/mol；deltaG1\_minus2\_350\_kJ = 0.25000000000000 kJ/mol；rate1\_increase\_factor = 31.410586383599 ratio。
\item 代回底物+产物的骨架守恒、概率归一或Eyring速率比，确认所有答案范围和符号。
\item 在长时间极限、没有反应极限及温度变化时，检查公式预言的物理趋势。
\item 本题所忽略的逆反应、构象重排、聚合动力学记忆和溶剂效应不应从理想数值随意外推。
\item 数值见证脚本验证输出一致性，独立实现和跨辑模型原创性仍需人工复核。
\end{enumerate}
\noindent\textbf{答案摘要}\par
k1\_298 = 11.7039954 s\textasciicircum{}-1；k2\_298 = 9.45967779 s\textasciicircum{}-1；k1\_350 = 367.62936 s\textasciicircum{}-1；k2\_350 = 400.608276 s\textasciicircum{}-1；branch\_ratio\_298 = 1.23725096 ratio；branch\_ratio\_350 = 0.917677897 ratio；P1\_fraction\_298 = 0.553022876 fraction；P1\_fraction\_350 = 0.478535993 fraction；deltaG1\_minus2\_298\_kJ = -0.52775 kJ/mol；deltaG1\_minus2\_350\_kJ = 0.25 kJ/mol；rate1\_increase\_factor = 31.4105864 ratio\par
\noindent\textbf{模型前提：}
本题无副反应和理想配分条件为假设；需独立科学审查。\par
\noindent\textcolor{red!65!black}{[跨辑去重未完成 · 草稿]} \par
\medskip\hrule\medskip
\Needspace{7\baselineskip}
\subsubsection*{UC-ADV15-O227 · 可逆烯胺形成的快速预平衡与有效亲核进攻速率}
\noindent\textbf{题目}\par
某酮K与催化胺形成烯胺E，二者快速预平衡E/K浓度比为q=0.5000（催化胺由储库维持恒定），随后只有烯胺E与Michael受体M发生不可逆二级加成，k2=4.000 L/(mol s)，自由M由外部库恒为0.01000 M。酮和烯胺合计底物浓度C0=0.02000 M，反应时间t=30.00 s，产物不逆转。求烯胺/酮分率、有效一级反应速度、底物剩余总浓度、产物浓度、30s转化率、产物初始生成速度及半衰期、如果预平衡比提升至q=2.000的kobs和速度倍率。强调快预平衡的底物中仅E可以反应，不能把k2[M]直接作用于全部K。\par
\medskip\noindent\textbf{分步解答}
\begin{enumerate}
\item 完整有机化学反应与初始组成：某酮K与催化胺形成烯胺E，二者快速预平衡E/K浓度比为q=0.5000（催化胺由储库维持恒定），随后只有烯胺E与Michael受体M发生不可逆二级加成，k2=4.000 L/(mol s)，自由M由外部库恒为0.01000 M。酮和烯胺合计底物浓度C0=0.02000 M，反应时间t=30.00 s，产物不逆转。求烯胺/酮分率、有效一级反应速度、底物剩余总浓度、产物浓度、30s转化率、产物初始生成速度及半衰期、如果预平衡比提升至q=2.000的kobs和速度倍率。强调快预平衡的底物中仅E可以反应，不能把k2[M]直接作用于全部K。
\item 从反应物库存、速率网络或统计假设建立数学式：fE=q/(1+q)，kobs=k2[M]fE，P=C0(1−e\textasciicircum{}(−kobs t))。
\item 写出底物与产物的严格物质守恒或反应构型概率，注意一级和二级速率的量纲转换。
\item 根据模型计算初始常数和分率：enamine\_fraction = 0.33333333333333 fraction；ketone\_fraction = 0.66666666666667 fraction；kobs\_s = 0.013333333333333 s\textasciicircum{}-1。
\item 继续计算反应进行后的浓度或分布：remaining\_total\_substrate\_M = 0.013406400920713 mol/L；product\_M = 0.0065935990792872 mol/L；conversion\_30 = 0.32967995396436 fraction。
\item 得到其余转化、选择性或分散度量：initial\_product\_rate\_M\_s = 0.00026666666666667 mol/(Ls)；half\_life\_s = 51.986038541996 s；kobs\_q2\_s = 0.026666666666667 s\textasciicircum{}-1；rate\_enhancement = 2.0000000000000 ratio。
\item 代回底物+产物的骨架守恒、概率归一或Eyring速率比，确认所有答案范围和符号。
\item 在长时间极限、没有反应极限及温度变化时，检查公式预言的物理趋势。
\item 本题所忽略的逆反应、构象重排、聚合动力学记忆和溶剂效应不应从理想数值随意外推。
\item 数值见证脚本验证输出一致性，独立实现和跨辑模型原创性仍需人工复核。
\end{enumerate}
\noindent\textbf{答案摘要}\par
enamine\_fraction = 0.333333333 fraction；ketone\_fraction = 0.666666667 fraction；kobs\_s = 0.0133333333 s\textasciicircum{}-1；remaining\_total\_substrate\_M = 0.0134064009 mol/L；product\_M = 0.00659359908 mol/L；conversion\_30 = 0.329679954 fraction；initial\_product\_rate\_M\_s = 0.000266666667 mol/(Ls)；half\_life\_s = 51.9860385 s；kobs\_q2\_s = 0.0266666667 s\textasciicircum{}-1；rate\_enhancement = 2 ratio\par
\noindent\textbf{模型前提：}
本题无副反应和理想配分条件为假设；需独立科学审查。\par
\noindent\textcolor{red!65!black}{[跨辑去重未完成 · 草稿]} \par
\medskip\hrule\medskip
\Needspace{7\baselineskip}
\subsubsection*{UC-ADV15-O228 · 具有二项序列分布的A/B共聚物三联体与序列熵}
\noindent\textbf{题目}\par
某理想A/B随机共聚物链节彼此独立，A出现概率p=0.4000，B概率q=0.6000，链足够长忽略端基。按独立Bernoulli序列模型，计算AA、AB、BA、BB相邻双联体概率；三联体ABA和BAB的概率、AAB及BAA合并三联体分率、每链节的二元Shannon序列熵h=−p ln p−q ln q（单位nat/链节）、长度N=1000理想样本中A链节的期望数和方差，并核对二联体概率总和。该模型无序列反应性比或链端记忆，真实自由基共聚物未必符合独立模型。\par
\medskip\noindent\textbf{分步解答}
\begin{enumerate}
\item 完整有机化学反应与初始组成：某理想A/B随机共聚物链节彼此独立，A出现概率p=0.4000，B概率q=0.6000，链足够长忽略端基。按独立Bernoulli序列模型，计算AA、AB、BA、BB相邻双联体概率；三联体ABA和BAB的概率、AAB及BAA合并三联体分率、每链节的二元Shannon序列熵h=−p ln p−q ln q（单位nat/链节）、长度N=1000理想样本中A链节的期望数和方差，并核对二联体概率总和。该模型无序列反应性比或链端记忆，真实自由基共聚物未必符合独立模型。
\item 从反应物库存、速率网络或统计假设建立数学式：P(AA)=p²，P(AB)=P(BA)=pq，P(BB)=q²，h=−p ln p−q ln q。
\item 写出底物与产物的严格物质守恒或反应构型概率，注意一级和二级速率的量纲转换。
\item 根据模型计算初始常数和分率：AA\_probability = 0.16000000000000 fraction；AB\_probability = 0.24000000000000 fraction；BA\_probability = 0.24000000000000 fraction。
\item 继续计算反应进行后的浓度或分布：BB\_probability = 0.36000000000000 fraction；ABA\_probability = 0.096000000000000 fraction；BAB\_probability = 0.14400000000000 fraction。
\item 得到其余转化、选择性或分散度量：AAB\_plus\_BAA\_probability = 0.19200000000000 fraction；entropy\_nat\_per\_monomer = 0.67301166700926 nat；expected\_A\_units\_N1000 = 400.00000000000 units；variance\_A\_units\_N1000 = 240.00000000000 units2；diad\_normalization = 1.0000000000000 ratio。
\item 代回底物+产物的骨架守恒、概率归一或Eyring速率比，确认所有答案范围和符号。
\item 在长时间极限、没有反应极限及温度变化时，检查公式预言的物理趋势。
\item 本题所忽略的逆反应、构象重排、聚合动力学记忆和溶剂效应不应从理想数值随意外推。
\item 数值见证脚本验证输出一致性，独立实现和跨辑模型原创性仍需人工复核。
\end{enumerate}
\noindent\textbf{答案摘要}\par
AA\_probability = 0.16 fraction；AB\_probability = 0.24 fraction；BA\_probability = 0.24 fraction；BB\_probability = 0.36 fraction；ABA\_probability = 0.096 fraction；BAB\_probability = 0.144 fraction；AAB\_plus\_BAA\_probability = 0.192 fraction；entropy\_nat\_per\_monomer = 0.673011667 nat；expected\_A\_units\_N1000 = 400 units；variance\_A\_units\_N1000 = 240 units2；diad\_normalization = 1 ratio\par
\noindent\textbf{模型前提：}
本题无副反应和理想配分条件为假设；需独立科学审查。\par
\noindent\textcolor{red!65!black}{[跨辑去重未完成 · 草稿]} \par
\medskip\hrule\medskip
\Needspace{7\baselineskip}
\subsubsection*{UC-ADV15-O229 · 随机断链有机聚合物的Poisson键断裂与剩余连通度}
\noindent\textbf{题目}\par
某线型有机聚合物原始每条链含N=500个重复单元和N−1=499个可断裂主链键。在给定照射剂量下，每个键独立发生不可逆断裂的概率p=0.01000，不形成交联与再接合。求每条原始链平均断键数、断键数方差、没有任何断键的链分率、至少一次断键的链分率、每条链平均产生的碎片链数E[fragments]=1+(N−1)p、某特定两个相邻键都完好的概率、任意选取的一条原始链主链键完好数量期望值，以及用Poisson近似λ=(N−1)p表示的零断键概率及与精确二项值的绝对差。说明零断裂的概率不等于单键存活概率。\par
\medskip\noindent\textbf{分步解答}
\begin{enumerate}
\item 完整有机化学反应与初始组成：某线型有机聚合物原始每条链含N=500个重复单元和N−1=499个可断裂主链键。在给定照射剂量下，每个键独立发生不可逆断裂的概率p=0.01000，不形成交联与再接合。求每条原始链平均断键数、断键数方差、没有任何断键的链分率、至少一次断键的链分率、每条链平均产生的碎片链数E[fragments]=1+(N−1)p、某特定两个相邻键都完好的概率、任意选取的一条原始链主链键完好数量期望值，以及用Poisson近似λ=(N−1)p表示的零断键概率及与精确二项值的绝对差。说明零断裂的概率不等于单键存活概率。
\item 从反应物库存、速率网络或统计假设建立数学式：X\textasciitilde{}Binomial(N−1,p)，P(X=0)=(1−p)\textasciicircum{}(N−1)，E[fragments]=1+E[X]。
\item 写出底物与产物的严格物质守恒或反应构型概率，注意一级和二级速率的量纲转换。
\item 根据模型计算初始常数和分率：mean\_breaks = 4.9900000000000 breaks；variance\_breaks = 4.9401000000000 breaks2；fraction\_intact\_chains = 0.0066368515579945 fraction。
\item 继续计算反应进行后的浓度或分布：fraction\_broken\_chains = 0.99336314844201 fraction；expected\_fragments = 5.9900000000000 chains；two\_adjacent\_bonds\_intact = 0.98010000000000 fraction。
\item 得到其余转化、选择性或分散度量：mean\_intact\_bonds = 494.01000000000 bonds；Poisson\_lambda = 4.9900000000000 breaks；Poisson\_P0 = 0.0068056644922305 fraction；exact\_minus\_poisson\_P0 = -0.00016881293423599 fraction。
\item 代回底物+产物的骨架守恒、概率归一或Eyring速率比，确认所有答案范围和符号。
\item 在长时间极限、没有反应极限及温度变化时，检查公式预言的物理趋势。
\item 本题所忽略的逆反应、构象重排、聚合动力学记忆和溶剂效应不应从理想数值随意外推。
\item 数值见证脚本验证输出一致性，独立实现和跨辑模型原创性仍需人工复核。
\end{enumerate}
\noindent\textbf{答案摘要}\par
mean\_breaks = 4.99 breaks；variance\_breaks = 4.9401 breaks2；fraction\_intact\_chains = 0.00663685156 fraction；fraction\_broken\_chains = 0.993363148 fraction；expected\_fragments = 5.99 chains；two\_adjacent\_bonds\_intact = 0.9801 fraction；mean\_intact\_bonds = 494.01 bonds；Poisson\_lambda = 4.99 breaks；Poisson\_P0 = 0.00680566449 fraction；exact\_minus\_poisson\_P0 = -0.000168812934 fraction\par
\noindent\textbf{模型前提：}
本题无副反应和理想配分条件为假设；需独立科学审查。\par
\noindent\textcolor{red!65!black}{[跨辑去重未完成 · 草稿]} \par
\medskip\hrule\medskip
\Needspace{7\baselineskip}
\subsubsection*{UC-ADV15-O230 · β位H/D氘代反应的并行KIE与产物同位素分流}
\noindent\textbf{题目}\par
某碱促E2消除反应中同结构H和D氘代底物的本征二级速度常数分别kH=5.000和kD=2.000 L/(mol s)。反应物H、D初始各0.01000 mol/L，碱由外部库恒定[B]=0.02000 M；H/D两底物平行独立且无同位素交换、没有反应后消除产物的逆反应。求H/D拟一级速度、反应t=20.00 s时两原料残余浓度、对应产物摩尔浓度、两通道转化率、H产物/D产物比、原始KIE速度比及t=0时两通道初始产物生成速率比。指出最终产物比会随时间变化，不能简单等同于瞬时KIE。\par
\medskip\noindent\textbf{分步解答}
\begin{enumerate}
\item 完整有机化学反应与初始组成：某碱促E2消除反应中同结构H和D氘代底物的本征二级速度常数分别kH=5.000和kD=2.000 L/(mol s)。反应物H、D初始各0.01000 mol/L，碱由外部库恒定[B]=0.02000 M；H/D两底物平行独立且无同位素交换、没有反应后消除产物的逆反应。求H/D拟一级速度、反应t=20.00 s时两原料残余浓度、对应产物摩尔浓度、两通道转化率、H产物/D产物比、原始KIE速度比及t=0时两通道初始产物生成速率比。指出最终产物比会随时间变化，不能简单等同于瞬时KIE。
\item 从反应物库存、速率网络或统计假设建立数学式：kobs,H=kH[B]，kobs,D=kD[B]，Pi(t)=C0(1−exp(−kobsi t))。
\item 写出底物与产物的严格物质守恒或反应构型概率，注意一级和二级速率的量纲转换。
\item 根据模型计算初始常数和分率：kobs\_H = 0.10000000000000 s\textasciicircum{}-1；kobs\_D = 0.040000000000000 s\textasciicircum{}-1；remaining\_H\_t20 = 0.0013533528323661 mol/L。
\item 继续计算反应进行后的浓度或分布：remaining\_D\_t20 = 0.0044932896411722 mol/L；product\_H\_t20 = 0.0086466471676339 mol/L；product\_D\_t20 = 0.0055067103588278 mol/L。
\item 得到其余转化、选择性或分散度量：conversion\_H = 0.86466471676339 fraction；conversion\_D = 0.55067103588278 fraction；product\_H\_over\_D = 1.5702019180603 ratio；intrinsic\_KIE = 2.5000000000000 ratio；initial\_rate\_ratio = 2.5000000000000 ratio。
\item 代回底物+产物的骨架守恒、概率归一或Eyring速率比，确认所有答案范围和符号。
\item 在长时间极限、没有反应极限及温度变化时，检查公式预言的物理趋势。
\item 本题所忽略的逆反应、构象重排、聚合动力学记忆和溶剂效应不应从理想数值随意外推。
\item 数值见证脚本验证输出一致性，独立实现和跨辑模型原创性仍需人工复核。
\end{enumerate}
\noindent\textbf{答案摘要}\par
kobs\_H = 0.1 s\textasciicircum{}-1；kobs\_D = 0.04 s\textasciicircum{}-1；remaining\_H\_t20 = 0.00135335283 mol/L；remaining\_D\_t20 = 0.00449328964 mol/L；product\_H\_t20 = 0.00864664717 mol/L；product\_D\_t20 = 0.00550671036 mol/L；conversion\_H = 0.864664717 fraction；conversion\_D = 0.550671036 fraction；product\_H\_over\_D = 1.57020192 ratio；intrinsic\_KIE = 2.5 ratio；initial\_rate\_ratio = 2.5 ratio\par
\noindent\textbf{模型前提：}
本题无副反应和理想配分条件为假设；需独立科学审查。\par
\noindent\textcolor{red!65!black}{[跨辑去重未完成 · 草稿]} \par
\medskip\hrule\medskip
\subsection*{物理化学}
\addcontentsline{toc}{subsection}{第15辑 / 物理化学}
\Needspace{7\baselineskip}
\subsubsection*{UC-ADV15-P223 · Gibbs–Helmholtz积分法对可逆络合反应平衡常数的温度校正}
\noindent\textbf{题目}\par
某均相络合反应A+L⇌AL在298.15 K标准热力学平衡常数K₁=100.0，以各溶解组分活动度除以1 M标准态定义K。若反应焓ΔH°=−20.00 kJ/mol在298–350 K范围恒定，气体常数R=8.314462618 J/(mol K)，求318.15与348.15 K的平衡常数、相对于298.15 K的K比、标准反应Gibbs自由能ΔG°298、ΔG°318、ΔG°348、350 K临近温区升温对络合平衡的影响方向、ln(K318/K298)及lnK对1/T的斜率−ΔH°/R。必须说明放热反应随温度上升K减小并且K无量纲。\par
\medskip\noindent\textbf{分步解答}
\begin{enumerate}
\item 给定热力学与物理化学模型：某均相络合反应A+L⇌AL在298.15 K标准热力学平衡常数K₁=100.0，以各溶解组分活动度除以1 M标准态定义K。若反应焓ΔH°=−20.00 kJ/mol在298–350 K范围恒定，气体常数R=8.314462618 J/(mol K)，求318.15与348.15 K的平衡常数、相对于298.15 K的K比、标准反应Gibbs自由能ΔG°298、ΔG°318、ΔG°348、350 K临近温区升温对络合平衡的影响方向、ln(K318/K298)及lnK对1/T的斜率−ΔH°/R。必须说明放热反应随温度上升K减小并且K无量纲。
\item 独立变量、近似条件与主方程：ln[K(T)/K(T1)]=−ΔH°/R·(1/T−1/T1)，ΔG°(T)=−RTlnK(T)。
\item 根据统计物理、动力学稳态或能量与物料守恒关系分离可观测量，检查单位转换和边界条件。
\item 计算基本状态参数：K\_298 = 100.00000000000 ratio；K\_318 = 60.219351425523 ratio；K\_348 = 31.389934211400 ratio；K318\_over\_K298 = 0.60219351425523 ratio。
\item 进一步得到动力学或热力学中间量：K348\_over\_K298 = 0.31389934211400 ratio；DG298\_kJmol = -11.416019004860 kJ/mol；DG318\_kJmol = -10.840202738206 kJ/mol。
\item 其余全部答案和回代量：DG348\_kJmol = -9.9764783382259 kJ/mol；log\_ratio\_318 = -0.50717643307681 log；vanthoff\_slope\_K = 2405.4471008989 K。
\item 利用统计分率归一、物质库存、能量差和热力学恒等式回代验证模型。
\item 考虑稀释、零反应速率、长时间或高低温极限，以发现符号、范围或极限不一致。
\item 真正的实验体系可能涉及非理想活动度、非稳态输运、电子耦合、结构重排或热容变化，题面近似不能越界。
\item 数值QA目前校验公式生成的完整结果，不将同源算法标注为完成独立实现；科学性及跨辑原创性未获专家批准。
\end{enumerate}
\noindent\textbf{答案摘要}\par
K\_298 = 100 ratio；K\_318 = 60.2193514 ratio；K\_348 = 31.3899342 ratio；K318\_over\_K298 = 0.602193514 ratio；K348\_over\_K298 = 0.313899342 ratio；DG298\_kJmol = -11.416019 kJ/mol；DG318\_kJmol = -10.8402027 kJ/mol；DG348\_kJmol = -9.97647834 kJ/mol；log\_ratio\_318 = -0.507176433 log；vanthoff\_slope\_K = 2405.4471 K\par
\noindent\textbf{模型前提：}
题设的理想相互作用和边界条件；专家科学与学术独立性审核未完成。\par
\noindent\textcolor{red!65!black}{[跨辑去重未完成 · 草稿]} \par
\medskip\hrule\medskip
\Needspace{7\baselineskip}
\subsubsection*{UC-ADV15-P224 · Einstein固体振动热容的量子高低温比较}
\noindent\textbf{题目}\par
某理想三维Einstein固体每mol原子具有3个相同频率的独立谐振子，Einstein特征温度ΘE=300.0 K。忽略电子热容与热膨胀，正则系综的摩尔定容热容C\_V=3R x² e\textasciicircum{}x/(e\textasciicircum{}x−1)²，x=ΘE/T。比较T=150.0和600.0 K，求各x、各量子热容、与经典Dulong–Petit极限3R的比值、两温度的单模平均激发态声子数n=1/(e\textasciicircum{}x−1)、两温度声子内能Uthermal=3RΘE n，以及热容高低温比。零点能不贡献C\_V但需在解释U时说明被扣除。\par
\medskip\noindent\textbf{分步解答}
\begin{enumerate}
\item 给定热力学与物理化学模型：某理想三维Einstein固体每mol原子具有3个相同频率的独立谐振子，Einstein特征温度ΘE=300.0 K。忽略电子热容与热膨胀，正则系综的摩尔定容热容C\_V=3R x² e\textasciicircum{}x/(e\textasciicircum{}x−1)²，x=ΘE/T。比较T=150.0和600.0 K，求各x、各量子热容、与经典Dulong–Petit极限3R的比值、两温度的单模平均激发态声子数n=1/(e\textasciicircum{}x−1)、两温度声子内能Uthermal=3RΘE n，以及热容高低温比。零点能不贡献C\_V但需在解释U时说明被扣除。
\item 独立变量、近似条件与主方程：Cv=3R x²exp(x)/(exp(x)−1)²，Uthermal=3RΘE/(exp(x)−1)。
\item 根据统计物理、动力学稳态或能量与物料守恒关系分离可观测量，检查单位转换和边界条件。
\item 计算基本状态参数：x\_150 = 2.0000000000000 ratio；x\_600 = 0.50000000000000 ratio；Cv\_150\_JmolK = 18.060550839694 J/(molK)；Cv\_600\_JmolK = 24.430165732405 J/(molK)。
\item 进一步得到动力学或热力学中间量：Cv150\_over\_classic = 0.72406166096631 ratio；Cv600\_over\_classic = 0.97942452225819 ratio；occupation\_150 = 0.15651764274967 phonons；occupation\_600 = 1.5414940825368 phonons。
\item 其余全部答案和回代量：Uthermal\_150\_Jmol = 1171.2240807296 J/mol；Uthermal\_600\_Jmol = 11535.025432608 J/mol；Cv\_high\_low\_ratio = 1.3526810975616 ratio。
\item 利用统计分率归一、物质库存、能量差和热力学恒等式回代验证模型。
\item 考虑稀释、零反应速率、长时间或高低温极限，以发现符号、范围或极限不一致。
\item 真正的实验体系可能涉及非理想活动度、非稳态输运、电子耦合、结构重排或热容变化，题面近似不能越界。
\item 数值QA目前校验公式生成的完整结果，不将同源算法标注为完成独立实现；科学性及跨辑原创性未获专家批准。
\end{enumerate}
\noindent\textbf{答案摘要}\par
x\_150 = 2 ratio；x\_600 = 0.5 ratio；Cv\_150\_JmolK = 18.0605508 J/(molK)；Cv\_600\_JmolK = 24.4301657 J/(molK)；Cv150\_over\_classic = 0.724061661 ratio；Cv600\_over\_classic = 0.979424522 ratio；occupation\_150 = 0.156517643 phonons；occupation\_600 = 1.54149408 phonons；Uthermal\_150\_Jmol = 1171.22408 J/mol；Uthermal\_600\_Jmol = 11535.0254 J/mol；Cv\_high\_low\_ratio = 1.3526811 ratio\par
\noindent\textbf{模型前提：}
题设的理想相互作用和边界条件；专家科学与学术独立性审核未完成。\par
\noindent\textcolor{red!65!black}{[跨辑去重未完成 · 草稿]} \par
\medskip\hrule\medskip
\Needspace{7\baselineskip}
\subsubsection*{UC-ADV15-P225 · 连续搅拌槽一级反应器的停留时间分布与稳态转化率}
\noindent\textbf{题目}\par
恒温理想连续搅拌反应器(CSTR)体积V=1.000 L，体积流量Q=0.05000 L/s，入口反应物浓度CA,in=0.1000 M，仅发生不可逆一级反应A→P，速率常数k=0.02000 s⁻¹。假设完全混合、出口浓度等于槽内浓度、体积与温度恒定。求停留时间τ=V/Q、Damköhler数Da=kτ、稳态出口A浓度与产物P浓度、反应转化率X=Da/(1+Da)、入口和出口A的摩尔通量(mol/s)、槽内反应物库存以及反应消耗摩尔速率kVCA和将同条件改为理想活塞流PFR时的转化率1−exp(−Da)。说明CSTR与PFR不同。\par
\medskip\noindent\textbf{分步解答}
\begin{enumerate}
\item 给定热力学与物理化学模型：恒温理想连续搅拌反应器(CSTR)体积V=1.000 L，体积流量Q=0.05000 L/s，入口反应物浓度CA,in=0.1000 M，仅发生不可逆一级反应A→P，速率常数k=0.02000 s⁻¹。假设完全混合、出口浓度等于槽内浓度、体积与温度恒定。求停留时间τ=V/Q、Damköhler数Da=kτ、稳态出口A浓度与产物P浓度、反应转化率X=Da/(1+Da)、入口和出口A的摩尔通量(mol/s)、槽内反应物库存以及反应消耗摩尔速率kVCA和将同条件改为理想活塞流PFR时的转化率1−exp(−Da)。说明CSTR与PFR不同。
\item 独立变量、近似条件与主方程：C\_A,out=C\_A,in/(1+kV/Q)，P\_out=C\_A,in−C\_A,out；rA=kVC\_A,out。
\item 根据统计物理、动力学稳态或能量与物料守恒关系分离可观测量，检查单位转换和边界条件。
\item 计算基本状态参数：residence\_time\_s = 20.000000000000 s；Damkohler\_number = 0.40000000000000 ratio；outlet\_A\_M = 0.071428571428571 mol/L；outlet\_P\_M = 0.028571428571429 mol/L。
\item 进一步得到动力学或热力学中间量：CSTR\_conversion = 0.28571428571429 fraction；inlet\_A\_flow\_mol\_s = 0.0050000000000000 mol/s；outlet\_A\_flow\_mol\_s = 0.0035714285714286 mol/s。
\item 其余全部答案和回代量：reactor\_A\_inventory\_mol = 0.071428571428571 mol；reaction\_rate\_mol\_s = 0.0014285714285714 mol/s；PFR\_conversion = 0.32967995396436 fraction。
\item 利用统计分率归一、物质库存、能量差和热力学恒等式回代验证模型。
\item 考虑稀释、零反应速率、长时间或高低温极限，以发现符号、范围或极限不一致。
\item 真正的实验体系可能涉及非理想活动度、非稳态输运、电子耦合、结构重排或热容变化，题面近似不能越界。
\item 数值QA目前校验公式生成的完整结果，不将同源算法标注为完成独立实现；科学性及跨辑原创性未获专家批准。
\end{enumerate}
\noindent\textbf{答案摘要}\par
residence\_time\_s = 20 s；Damkohler\_number = 0.4 ratio；outlet\_A\_M = 0.0714285714 mol/L；outlet\_P\_M = 0.0285714286 mol/L；CSTR\_conversion = 0.285714286 fraction；inlet\_A\_flow\_mol\_s = 0.005 mol/s；outlet\_A\_flow\_mol\_s = 0.00357142857 mol/s；reactor\_A\_inventory\_mol = 0.0714285714 mol；reaction\_rate\_mol\_s = 0.00142857143 mol/s；PFR\_conversion = 0.329679954 fraction\par
\noindent\textbf{模型前提：}
题设的理想相互作用和边界条件；专家科学与学术独立性审核未完成。\par
\noindent\textcolor{red!65!black}{[跨辑去重未完成 · 草稿]} \par
\medskip\hrule\medskip
\Needspace{7\baselineskip}
\subsubsection*{UC-ADV15-P226 · Marcus电子转移动力学的正常区、活化零点与反转区}
\noindent\textbf{题目}\par
一电子非绝热电子转移反应的Marcus经典活化自由能ΔG‡=(λ+ΔG°)²/(4λ)，重组能λ=0.8000 eV。取三种驱动力ΔG°分别−0.2000、−0.8000和−1.400 eV，在T=298.15 K假设相同的电子耦合与频率前因子k0=1.000×10¹¹ s⁻¹，使用k=k0 exp(−ΔG‡/(kBT))，kB=8.617333262145×10⁻⁵ eV/K。求三种驱动力下活化自由能、速率、−0.8 eV状态的零势垒极值速率、正常区与反转区的速率比、常温kBT(eV)以及相对于最优速率的无量纲两状态速率。说明该模型的反转区反映势垒重新增加，并非违反放热自发性。\par
\medskip\noindent\textbf{分步解答}
\begin{enumerate}
\item 给定热力学与物理化学模型：一电子非绝热电子转移反应的Marcus经典活化自由能ΔG‡=(λ+ΔG°)²/(4λ)，重组能λ=0.8000 eV。取三种驱动力ΔG°分别−0.2000、−0.8000和−1.400 eV，在T=298.15 K假设相同的电子耦合与频率前因子k0=1.000×10¹¹ s⁻¹，使用k=k0 exp(−ΔG‡/(kBT))，kB=8.617333262145×10⁻⁵ eV/K。求三种驱动力下活化自由能、速率、−0.8 eV状态的零势垒极值速率、正常区与反转区的速率比、常温kBT(eV)以及相对于最优速率的无量纲两状态速率。说明该模型的反转区反映势垒重新增加，并非违反放热自发性。
\item 独立变量、近似条件与主方程：ΔG‡=(λ+ΔG°)²/(4λ)；k=k0exp(−ΔG‡/kBT)。
\item 根据统计物理、动力学稳态或能量与物料守恒关系分离可观测量，检查单位转换和边界条件。
\item 计算基本状态参数：activation\_normal\_eV = 0.11250000000000 eV；activation\_optimum\_eV = 0.0000000000000 eV；activation\_inverted\_eV = 0.11250000000000 eV；rate\_normal\_s = 1254169913.4018 s\textasciicircum{}-1。
\item 进一步得到动力学或热力学中间量：rate\_optimum\_s = 100000000000.00 s\textasciicircum{}-1；rate\_inverted\_s = 1254169913.4018 s\textasciicircum{}-1；rate\_normal\_over\_inverted = 1.0000000000000 ratio。
\item 其余全部答案和回代量：kBT\_eV = 0.025692579121085 eV；normal\_relative\_to\_optimum = 0.012541699134018 ratio；inverted\_relative\_to\_optimum = 0.012541699134018 ratio。
\item 利用统计分率归一、物质库存、能量差和热力学恒等式回代验证模型。
\item 考虑稀释、零反应速率、长时间或高低温极限，以发现符号、范围或极限不一致。
\item 真正的实验体系可能涉及非理想活动度、非稳态输运、电子耦合、结构重排或热容变化，题面近似不能越界。
\item 数值QA目前校验公式生成的完整结果，不将同源算法标注为完成独立实现；科学性及跨辑原创性未获专家批准。
\end{enumerate}
\noindent\textbf{答案摘要}\par
activation\_normal\_eV = 0.1125 eV；activation\_optimum\_eV = 0 eV；activation\_inverted\_eV = 0.1125 eV；rate\_normal\_s = 1.25416991e+09 s\textasciicircum{}-1；rate\_optimum\_s = 1e+11 s\textasciicircum{}-1；rate\_inverted\_s = 1.25416991e+09 s\textasciicircum{}-1；rate\_normal\_over\_inverted = 1 ratio；kBT\_eV = 0.0256925791 eV；normal\_relative\_to\_optimum = 0.0125416991 ratio；inverted\_relative\_to\_optimum = 0.0125416991 ratio\par
\noindent\textbf{模型前提：}
题设的理想相互作用和边界条件；专家科学与学术独立性审核未完成。\par
\noindent\textcolor{red!65!black}{[跨辑去重未完成 · 草稿]} \par
\medskip\hrule\medskip
\Needspace{7\baselineskip}
\subsubsection*{UC-ADV15-P227 · Lindemann–Hinshelwood单分子反应的低压极限和高压饱和}
\noindent\textbf{题目}\par
气相分子A的单分子异构化采用Lindemann机制A+M→A*+M激活，A*+M→A+M失活和A*→P转化。给定k1=1.000×10⁻¹² cm³/(molecule s)、k−1=2.000×10⁻¹² cm³/(molecule s)、k2=1.000×10⁴ s⁻¹，第三体M的数密度分别为n1=1.000×10¹⁵与n2=1.000×10¹⁷ molecule/cm³。准稳态条件下有效一级速率kobs=k1 n k2/(k−1 n+k2)。求两个n下的激活频率、失活频率、有效速度、对高压极限k∞=k1 k2/k−1的归一分率、低压线性近似k1 n的值、高低压力速率比及半饱和第三体密度n½=k2/k−1。不能把实验低压近似应用到高压。\par
\medskip\noindent\textbf{分步解答}
\begin{enumerate}
\item 给定热力学与物理化学模型：气相分子A的单分子异构化采用Lindemann机制A+M→A*+M激活，A*+M→A+M失活和A*→P转化。给定k1=1.000×10⁻¹² cm³/(molecule s)、k−1=2.000×10⁻¹² cm³/(molecule s)、k2=1.000×10⁴ s⁻¹，第三体M的数密度分别为n1=1.000×10¹⁵与n2=1.000×10¹⁷ molecule/cm³。准稳态条件下有效一级速率kobs=k1 n k2/(k−1 n+k2)。求两个n下的激活频率、失活频率、有效速度、对高压极限k∞=k1 k2/k−1的归一分率、低压线性近似k1 n的值、高低压力速率比及半饱和第三体密度n½=k2/k−1。不能把实验低压近似应用到高压。
\item 独立变量、近似条件与主方程：kobs=(k1 n)k2/(k−1 n+k2)，k∞=k1 k2/k−1。
\item 根据统计物理、动力学稳态或能量与物料守恒关系分离可观测量，检查单位转换和边界条件。
\item 计算基本状态参数：activation\_freq\_low = 1000.0000000000 s\textasciicircum{}-1；activation\_freq\_high = 100000.00000000 s\textasciicircum{}-1；deactivation\_freq\_low = 2000.0000000000 s\textasciicircum{}-1；deactivation\_freq\_high = 200000.00000000 s\textasciicircum{}-1。
\item 进一步得到动力学或热力学中间量：kobs\_low = 833.33333333333 s\textasciicircum{}-1；kobs\_high = 4761.9047619048 s\textasciicircum{}-1；high\_pressure\_limit = 5000.0000000000 s\textasciicircum{}-1；fraction\_of\_limit\_low = 0.16666666666667 fraction。
\item 其余全部答案和回代量：fraction\_of\_limit\_high = 0.95238095238095 fraction；pressure\_rate\_ratio = 5.7142857142857 ratio；half\_density = 5.0000000000000e+15 molecules/cm3。
\item 利用统计分率归一、物质库存、能量差和热力学恒等式回代验证模型。
\item 考虑稀释、零反应速率、长时间或高低温极限，以发现符号、范围或极限不一致。
\item 真正的实验体系可能涉及非理想活动度、非稳态输运、电子耦合、结构重排或热容变化，题面近似不能越界。
\item 数值QA目前校验公式生成的完整结果，不将同源算法标注为完成独立实现；科学性及跨辑原创性未获专家批准。
\end{enumerate}
\noindent\textbf{答案摘要}\par
activation\_freq\_low = 1000 s\textasciicircum{}-1；activation\_freq\_high = 100000 s\textasciicircum{}-1；deactivation\_freq\_low = 2000 s\textasciicircum{}-1；deactivation\_freq\_high = 200000 s\textasciicircum{}-1；kobs\_low = 833.333333 s\textasciicircum{}-1；kobs\_high = 4761.90476 s\textasciicircum{}-1；high\_pressure\_limit = 5000 s\textasciicircum{}-1；fraction\_of\_limit\_low = 0.166666667 fraction；fraction\_of\_limit\_high = 0.952380952 fraction；pressure\_rate\_ratio = 5.71428571 ratio；half\_density = 5e+15 molecules/cm3\par
\noindent\textbf{模型前提：}
题设的理想相互作用和边界条件；专家科学与学术独立性审核未完成。\par
\noindent\textcolor{red!65!black}{[跨辑去重未完成 · 草稿]} \par
\medskip\hrule\medskip
\Needspace{7\baselineskip}
\subsubsection*{UC-ADV15-P228 · 气体Henry定律与温度溶解度的van’t Hoff校正}
\noindent\textbf{题目}\par
某中性气体在稀溶液与理想气相间满足c=kH(T)p，c以mol/L、p以bar、kH以mol/(L bar)计。298.15 K时kH=0.001000 mol/(L bar)，气体溶解标准焓ΔHsol=−15.00 kJ/mol且在温区内恒定，假定dlnkH/d(1/T)=−ΔHsol/R。比较T1=298.15和T2=318.15 K、p=1.000及5.000 bar，求kH两温度、各压强四个溶解浓度、温度引起的溶解度比、固定p=5 bar时温升导致的溶解量差、由数值反推的lnkH比和Henry常数斜率。需清楚本题kH以c/p定义，不能误用相反的p/c约定。\par
\medskip\noindent\textbf{分步解答}
\begin{enumerate}
\item 给定热力学与物理化学模型：某中性气体在稀溶液与理想气相间满足c=kH(T)p，c以mol/L、p以bar、kH以mol/(L bar)计。298.15 K时kH=0.001000 mol/(L bar)，气体溶解标准焓ΔHsol=−15.00 kJ/mol且在温区内恒定，假定dlnkH/d(1/T)=−ΔHsol/R。比较T1=298.15和T2=318.15 K、p=1.000及5.000 bar，求kH两温度、各压强四个溶解浓度、温度引起的溶解度比、固定p=5 bar时温升导致的溶解量差、由数值反推的lnkH比和Henry常数斜率。需清楚本题kH以c/p定义，不能误用相反的p/c约定。
\item 独立变量、近似条件与主方程：ln[kH(T)/kH(T1)]=−ΔHsol/R (1/T−1/T1)，c=kH p。
\item 根据统计物理、动力学稳态或能量与物料守恒关系分离可观测量，检查单位转换和边界条件。
\item 计算基本状态参数：kH\_298 = 0.0010000000000000 mol/(Lbar)；kH\_318 = 0.00068360000200506 mol/(Lbar)；c\_298\_1bar = 0.0010000000000000 mol/L；c\_298\_5bar = 0.0050000000000000 mol/L。
\item 进一步得到动力学或热力学中间量：c\_318\_1bar = 0.00068360000200506 mol/L；c\_318\_5bar = 0.0034180000100253 mol/L；temperature\_solubility\_ratio = 0.68360000200506 ratio。
\item 其余全部答案和回代量：concentration\_drop\_5bar = 0.0015819999899747 mol/L；ln\_kH\_ratio = -0.38038232480761 log；inverseT\_slope\_K = 1804.0853256741 K。
\item 利用统计分率归一、物质库存、能量差和热力学恒等式回代验证模型。
\item 考虑稀释、零反应速率、长时间或高低温极限，以发现符号、范围或极限不一致。
\item 真正的实验体系可能涉及非理想活动度、非稳态输运、电子耦合、结构重排或热容变化，题面近似不能越界。
\item 数值QA目前校验公式生成的完整结果，不将同源算法标注为完成独立实现；科学性及跨辑原创性未获专家批准。
\end{enumerate}
\noindent\textbf{答案摘要}\par
kH\_298 = 0.001 mol/(Lbar)；kH\_318 = 0.000683600002 mol/(Lbar)；c\_298\_1bar = 0.001 mol/L；c\_298\_5bar = 0.005 mol/L；c\_318\_1bar = 0.000683600002 mol/L；c\_318\_5bar = 0.00341800001 mol/L；temperature\_solubility\_ratio = 0.683600002 ratio；concentration\_drop\_5bar = 0.00158199999 mol/L；ln\_kH\_ratio = -0.380382325 log；inverseT\_slope\_K = 1804.08533 K\par
\noindent\textbf{模型前提：}
题设的理想相互作用和边界条件；专家科学与学术独立性审核未完成。\par
\noindent\textcolor{red!65!black}{[跨辑去重未完成 · 草稿]} \par
\medskip\hrule\medskip
\Needspace{7\baselineskip}
\subsubsection*{UC-ADV15-P229 · 一维无限深势阱电子能级与光吸收波长}
\noindent\textbf{题目}\par
一个电子被限制在宽度L=1.000 nm的一维无限深方势阱中，粒子质量me=9.1093837139×10⁻³¹ kg，Planck常数h=6.62607015×10⁻³⁴ J·s。能级En=n²h²/(8me L²)，n=1、2、3；取1 eV=1.602176634×10⁻¹⁹ J、光速c=299792458 m/s。求E1、E2、E3(eV)、1→2和2→3跃迁能隙(eV)、1→2跃迁光子波长(nm)、1→3能隙与波长、基态的零点量子动能(J)以及E3/E1。需区别无限阱能级与真实分子HOMO-LUMO，不包括电子相互作用。\par
\medskip\noindent\textbf{分步解答}
\begin{enumerate}
\item 给定热力学与物理化学模型：一个电子被限制在宽度L=1.000 nm的一维无限深方势阱中，粒子质量me=9.1093837139×10⁻³¹ kg，Planck常数h=6.62607015×10⁻³⁴ J·s。能级En=n²h²/(8me L²)，n=1、2、3；取1 eV=1.602176634×10⁻¹⁹ J、光速c=299792458 m/s。求E1、E2、E3(eV)、1→2和2→3跃迁能隙(eV)、1→2跃迁光子波长(nm)、1→3能隙与波长、基态的零点量子动能(J)以及E3/E1。需区别无限阱能级与真实分子HOMO-LUMO，不包括电子相互作用。
\item 独立变量、近似条件与主方程：En=n²h²/(8meL²)，ΔEij=Ej−Ei，λph=hc/ΔE。
\item 根据统计物理、动力学稳态或能量与物料守恒关系分离可观测量，检查单位转换和边界条件。
\item 计算基本状态参数：E1\_eV = 0.37603016210487 eV；E2\_eV = 1.5041206484195 eV；E3\_eV = 3.3842714589439 eV；gap12\_eV = 1.1280904863146 eV。
\item 进一步得到动力学或热力学中间量：gap23\_eV = 1.8801508105244 eV；lambda12\_nm = 1099.0625303316 nm；gap13\_eV = 3.0082412968390 eV。
\item 其余全部答案和回代量：lambda13\_nm = 412.14844887437 nm；E1\_J = 6.0246673940366e-20 J；E3\_over\_E1 = 9.0000000000000 ratio。
\item 利用统计分率归一、物质库存、能量差和热力学恒等式回代验证模型。
\item 考虑稀释、零反应速率、长时间或高低温极限，以发现符号、范围或极限不一致。
\item 真正的实验体系可能涉及非理想活动度、非稳态输运、电子耦合、结构重排或热容变化，题面近似不能越界。
\item 数值QA目前校验公式生成的完整结果，不将同源算法标注为完成独立实现；科学性及跨辑原创性未获专家批准。
\end{enumerate}
\noindent\textbf{答案摘要}\par
E1\_eV = 0.376030162 eV；E2\_eV = 1.50412065 eV；E3\_eV = 3.38427146 eV；gap12\_eV = 1.12809049 eV；gap23\_eV = 1.88015081 eV；lambda12\_nm = 1099.06253 nm；gap13\_eV = 3.0082413 eV；lambda13\_nm = 412.148449 nm；E1\_J = 6.02466739e-20 J；E3\_over\_E1 = 9 ratio\par
\noindent\textbf{模型前提：}
题设的理想相互作用和边界条件；专家科学与学术独立性审核未完成。\par
\noindent\textcolor{red!65!black}{[跨辑去重未完成 · 草稿]} \par
\medskip\hrule\medskip
\Needspace{7\baselineskip}
\subsubsection*{UC-ADV15-P230 · 二元液体超额摩尔体积的部分摩尔体积计算}
\noindent\textbf{题目}\par
某298.15 K的不可压缩二元液体A/B混合物，其摩尔体积VM(x)=xVA*+(1−x)VB*+Kx(1−x)，纯A摩尔体积VA*=80.00 cm³/mol、纯B的VB*=100.0 cm³/mol，非理想超额体积系数K=+20.00 cm³/mol。对xA=0.3000与0.7000求混合摩尔体积、超额摩尔体积、导数dVM/dx，以及部分摩尔体积vbarA=VM+(1−x)dVM/dx、vbarB=VM−x dVM/dx。再在xA=0.3验算x vbarA+(1−x)vbarB=VM与超额体积正负号。说明部分摩尔体积不同于纯物质体积，题面仅为给定的多项式模型。\par
\medskip\noindent\textbf{分步解答}
\begin{enumerate}
\item 给定热力学与物理化学模型：某298.15 K的不可压缩二元液体A/B混合物，其摩尔体积VM(x)=xVA*+(1−x)VB*+Kx(1−x)，纯A摩尔体积VA*=80.00 cm³/mol、纯B的VB*=100.0 cm³/mol，非理想超额体积系数K=+20.00 cm³/mol。对xA=0.3000与0.7000求混合摩尔体积、超额摩尔体积、导数dVM/dx，以及部分摩尔体积vbarA=VM+(1−x)dVM/dx、vbarB=VM−x dVM/dx。再在xA=0.3验算x vbarA+(1−x)vbarB=VM与超额体积正负号。说明部分摩尔体积不同于纯物质体积，题面仅为给定的多项式模型。
\item 独立变量、近似条件与主方程：VM=xVA*+(1−x)VB*+Kx(1−x)，VM'=VA*−VB*+K(1−2x)，vbarA=VM+(1−x)VM'，vbarB=VM−xVM'。
\item 根据统计物理、动力学稳态或能量与物料守恒关系分离可观测量，检查单位转换和边界条件。
\item 计算基本状态参数：Vm\_x0p3 = 98.200000000000 cm3/mol；Vm\_x0p7 = 90.200000000000 cm3/mol；excess\_x0p3 = 4.2000000000000 cm3/mol；excess\_x0p7 = 4.2000000000000 cm3/mol。
\item 进一步得到动力学或热力学中间量：derivative\_x0p3 = -12.000000000000 cm3/mol；derivative\_x0p7 = -28.000000000000 cm3/mol；partial\_A\_x0p3 = 89.800000000000 cm3/mol；partial\_B\_x0p3 = 101.80000000000 cm3/mol。
\item 其余全部答案和回代量：partial\_A\_x0p7 = 81.800000000000 cm3/mol；partial\_B\_x0p7 = 109.80000000000 cm3/mol；Euler\_balance\_x0p3 = 98.200000000000 cm3/mol；positive\_excess\_flag = 1.0000000000000 boolean。
\item 利用统计分率归一、物质库存、能量差和热力学恒等式回代验证模型。
\item 考虑稀释、零反应速率、长时间或高低温极限，以发现符号、范围或极限不一致。
\item 真正的实验体系可能涉及非理想活动度、非稳态输运、电子耦合、结构重排或热容变化，题面近似不能越界。
\item 数值QA目前校验公式生成的完整结果，不将同源算法标注为完成独立实现；科学性及跨辑原创性未获专家批准。
\end{enumerate}
\noindent\textbf{答案摘要}\par
Vm\_x0p3 = 98.2 cm3/mol；Vm\_x0p7 = 90.2 cm3/mol；excess\_x0p3 = 4.2 cm3/mol；excess\_x0p7 = 4.2 cm3/mol；derivative\_x0p3 = -12 cm3/mol；derivative\_x0p7 = -28 cm3/mol；partial\_A\_x0p3 = 89.8 cm3/mol；partial\_B\_x0p3 = 101.8 cm3/mol；partial\_A\_x0p7 = 81.8 cm3/mol；partial\_B\_x0p7 = 109.8 cm3/mol；Euler\_balance\_x0p3 = 98.2 cm3/mol；positive\_excess\_flag = 1 boolean\par
\noindent\textbf{模型前提：}
题设的理想相互作用和边界条件；专家科学与学术独立性审核未完成。\par
\noindent\textcolor{red!65!black}{[跨辑去重未完成 · 草稿]} \par
\medskip\hrule\medskip
\end{document}
